%%%%%%%% ICML 2026 EXAMPLE LATEX SUBMISSION FILE %%%%%%%%%%%%%%%%%

\documentclass{article}

% Recommended, but optional, packages for figures and better typesetting:
\usepackage{microtype}
\usepackage{graphicx}
\usepackage{booktabs}
\usepackage{makecell}
\usepackage{wrapfig}

% hyperref makes hyperlinks in the resulting PDF.
% If your build breaks (sometimes temporarily if a hyperlink spans a page)
% please comment out the following usepackage line and replace
% \usepackage{icml2026} with \usepackage[nohyperref]{icml2026} above.
\usepackage{hyperref}

% Attempt to make hyperref and algorithmic work together better:
\newcommand{\theHalgorithm}{\arabic{algorithm}}

% Authored preprint. Use \usepackage{icml2026} for a blind submission.
\usepackage[preprint]{icml2026}
\newcommand{\icmlpreprinttext}{} 
\renewcommand{\icmlpreprinttext}{}

% If accepted, instead use the following line for the camera-ready submission:
% \usepackage[accepted]{icml2026}

\usepackage{amsmath}
\usepackage{amssymb}
\usepackage{mathtools}
\usepackage{amsthm}


% if you use cleveref..
\usepackage[capitalize,noabbrev]{cleveref}
\providecommand{\Description}[1]{}

%%%%%%%%%%%%%%%%%%%%%%%%%%%%%%%%
% THEOREMS
%%%%%%%%%%%%%%%%%%%%%%%%%%%%%%%%
\theoremstyle{plain}
\newtheorem{theorem}{Theorem}[section]
\newtheorem{proposition}[theorem]{Proposition}
\newtheorem{lemma}[theorem]{Lemma}
\newtheorem{corollary}[theorem]{Corollary}
\theoremstyle{definition}
\newtheorem{definition}[theorem]{Definition}
\newtheorem{assumption}[theorem]{Assumption}
\theoremstyle{remark}
\newtheorem{remark}[theorem]{Remark}

% The \icmltitle you define below is probably too long as a header.
% Therefore, a short form for the running title is supplied here:
\icmltitlerunning{Remoteness-Aware Control of Negative Off-Policy Updates}
\setlength{\abovedisplayskip}{3pt}
\setlength{\belowdisplayskip}{3pt}

\begin{document}

\twocolumn[
  \icmltitle{Breaking the Curse of Repulsion: Remoteness-Aware Control of Negative Off-Policy Updates}

  % It is OKAY to include author information, even for blind submissions: the
  % style file will automatically remove it for you unless you've provided
  % the [accepted] option to the icml2026 package.

  % List of affiliations: The first argument should be a (short) identifier you
  % will use later to specify author affiliations Academic affiliations
  % should list Department, University, City, Region, Country Industry
  % affiliations should list Company, City, Region, Country

  % You can specify symbols, otherwise they are numbered in order. Ideally, you
  % should not use this facility. Affiliations will be numbered in order of
  % appearance and this is the preferred way.
  \begin{icmlauthorlist}
    \icmlauthor{Yusen Huo}{comp}
    \icmlauthor{Changping Wang}{comp}
    \icmlauthor{Yangru Huang}{comp}
    \icmlauthor{Jun Zhang}{comp}
    \icmlauthor{Jie Jiang}{comp}
  \end{icmlauthorlist}
  \icmlaffiliation{comp}{Tencent Inc, China}

  \icmlcorrespondingauthor{Jun Zhang}{neoxzhang@tencent.com}

  % You may provide any keywords that you find helpful for describing your
  % paper; these are used to populate the "keywords" metadata in the PDF but
  % will not be shown in the document
  \icmlkeywords{Off-Policy Reinforcement Learning, Negative Feedback, Policy Optimization, Stability}

  \vskip 0.3in
]

% this must go after the closing bracket ] following \twocolumn[ ...

% This command actually creates the footnote in the first column listing the
% affiliations and the copyright notice. The command takes one argument, which
% is text to display at the start of the footnote. The \icmlEqualContribution
% command is standard text for equal contribution. Remove it (just {}) if you
% do not need this facility.

% Use ONE of the following lines. DO NOT remove the command.
% If you have no special notice, KEEP empty braces:
% \printAffiliationsAndNotice{}  % no special notice (required even if empty)
% Or, if applicable, use the standard equal contribution text:

\printAffiliationsAndNotice{}

\begin{abstract}
Off-policy policy optimization reuses historical behavior, including
negative-advantage samples that suppress known failures. We show that repeated
reuse can turn this useful signal into excessive repulsion: as the learner
moves away from a historical negative action, subsequent updates make that
action increasingly remote without necessarily reducing its update strength.
Our aggregate theory characterizes the resulting transition from a stable
displacement beyond the positive-only target to persistent drift and the loss
of finite stable equilibria; controlled strength sweeps show that an
intermediate displacement can improve held-out reward. The relevant
learner-relative coordinate is squared standardized distance for Gaussian
policies and surprisal for categorical policies. We introduce Dynamic
Remoteness-Aware Policy Optimization (DRPO), which leaves the negative update
unchanged in the near field and exponentially attenuates its remote tail.
DRPO restores eventual inward Gaussian drift for every fixed finite
negative-to-positive mass ratio, yields an explicit ultimate-bound radius, and
changes categorical support suppression from exponential to polynomial
probability decay. External diagnostics and controlled interventions isolate
remoteness-dependent policy geometry as a source of negative-update
amplification and show that selective tapering can remove its far-field effect
without discarding useful local feedback.
\end{abstract}

\section{Introduction}
\label{sec:introduction}

Policy-based methods learn from signed feedback: positive updates attract the
policy toward actions estimated to be preferable, whereas negative updates
suppress actions estimated to be undesirable. Negative feedback can therefore encode useful boundaries and exclusion
information that positive-only learning lacks
\citep{zhu2025negative,song2025good}. Such information can displace the
policy beyond the solution supported by positive feedback alone, although
whether that displacement improves task utility is an empirical question. In off-policy settings, however, the same historical sample may be reused
after earlier updates have already changed its probability under the current
policy; we study the zero-refresh limit for controlled isolation. This raises a single
question: when does repeated reuse turn useful negative feedback into
excessive repulsion?

We identify learner-relative remoteness as the state variable that links one
reuse step to the next. A negative update lowers the probability of its historical action,
increasing the action's remoteness and thereby changing the next
policy-score response (the \emph{near field} \(D\leq\tau\) versus the
\emph{far field} \(D>\tau\); Section~\ref{sec:drpo}). The policy family determines the resulting tail
law: Gaussian mean scores grow with standardized distance, whereas categorical
logit scores remain bounded even as selected-action probabilities approach the
simplex boundary. Under isolated reuse, uncontrolled Gaussian remoteness grows
geometrically and categorical action probability decays exponentially; an
exponential taper reduces the corresponding far-field growth to logarithmic
order. This is the curse of repulsion: following a negative update makes its next
reuse more remote without making the signal self-limiting.

What remains unresolved is whether remoteness itself amplifies negative
updates independently of sample quality. Existing methods control off-policy
instability through filtering, importance weighting, probability-aware gates,
and tapered updates
\citep{novati2019remember,fakoor2020p3o,peng2019advantage,
schulman2017proximal,fujimoto2019off,kumar2020conservative,
arnal2025asymmetric,leroux2025tapered}, but remoteness is usually confounded
with advantage magnitude and data quality. We therefore match those factors
while varying only remoteness, then intervene separately on near- and far-field
negative updates to test their downstream effects.

Our aggregate analysis places these per-sample effects in competition with
positive attraction. When positive mass dominates, repulsion can displace the
policy beyond the positive-only target while preserving a finite stable
equilibrium. At balance, the restoring force disappears and persistent drift
emerges; under negative dominance, no finite stable equilibrium remains.
These regimes characterize policy geometry and stability, not task utility,
which must be determined empirically.

Motivated by this transition, we introduce Dynamic Remoteness-Aware Policy
Optimization (DRPO). DRPO recomputes remoteness under the current policy,
leaves the negative branch unchanged in the near field, and applies an
exponential taper beyond the threshold. The taper dominates every finite-order
far-field response, restores eventual inward Gaussian drift for any fixed
finite negative-to-positive mass ratio, and makes categorical suppression
self-limiting. Remoteness controls geometric exposure, not directional utility
(Section~\ref{subsec:distance_dependent_utility}).

We first verify the pattern externally, then use controlled environments to
isolate its source, test its downstream effect, and compare control rules.

Our contributions are: \textbf{(1)} we characterize how repeated negative
reuse changes policy-relative remoteness and derive aggregate regimes for
stable extrapolation, persistent drift, and the loss of finite stable
equilibria; \textbf{(2)} we derive a control-order hierarchy and a selective
exponential taper with Gaussian ultimate-boundedness and categorical
self-limiting suppression guarantees; and \textbf{(3)} matched external
diagnostics and controlled interventions separate coefficient magnitude
from policy geometry, identify a far-field causal pathway, and test the
stability--utility trade-off of selective tapering.


% ------------------------------------------------------------------
% SECTION 2: PRELIMINARIES
% ------------------------------------------------------------------
\section{Related Work}
\label{sec:related_work}

\textbf{Off-policy historical reuse across policy domains.}
Off-policy policy optimization reuses data generated by earlier policies,
from replay and offline-RL datasets in continuous control to fixed
preference and verifier-labeled corpora in language-model post-training.
Classical approaches correct behavior--learner mismatch or constrain policy
support through importance weighting, clipping, trust regions, and
conservative or behavior-regularized objectives
\citep{precup2000eligibility,schulman2015trust,
schulman2017proximal,fujimoto2019off,kumar2020conservative,
fujimoto2021minimalist,kostrikov2021offline}. Truncated importance ratios also underlie corrected off-policy return
estimators such as Retrace and V-trace
\citep{munos2016safe,espeholt2018impala}; their primary role is
distribution correction rather than learner-relative negative-tail control. Across both continuous and
discrete policies, our focus is narrower: repeated reuse of signed actor
feedback, especially the negative branch, after the learner has moved away
from the historical action.

\textbf{Controlling negative updates.}
Existing controls span positive filtering, which removes the negative branch
\citep{peng2019advantage,nair2020awac,kostrikov2021offline}; fixed global
down-weighting of negative advantages, as in HPO's hysteretic weight
\citep{sana2026hpo}; baseline-based positive--negative rebalancing in
AsymRE \citep{arnal2025asymmetric}; reference-relative preference
optimization in DPO \citep{rafailov2023direct}; and probability-aware
tapering in TOPR \citep{leroux2025tapered}. Failed or rejected behavior is
also exploited in unlikelihood and preference-based training
\citep{welleck2019neural,ethayarajh2024kto,duan2024negating,
zhang2024negative,zhu2025negative,song2025good}.
DRPO differs by treating learner-relative remoteness as a dynamical state
created by repeated reuse and by selectively attenuating only the far-field
negative tail.
% ------------------------------------------------------------------
% SECTION 2: PRELIMINARIES
% ------------------------------------------------------------------

\section{Problem Setup}
\label{sec:problem_setup}

% ------------------------------------------------------------------
% SECTION 2: PRELIMINARIES
% ------------------------------------------------------------------

\subsection{Off-Policy Signed Actor Updates}

We consider policy optimization from a historical update distribution
$\nu$ over state--action pairs, which may represent an offline dataset,
a replay buffer, or trajectories collected by earlier or stale behavior
policies. Let $\pi_\theta(a \mid s)$ denote the current policy and let
$\widehat{A}(s,a)$ denote the advantage-like signal used by the actor
update. During an actor step, the update distribution and the associated
advantage estimates are treated as fixed with respect to $\theta$. The resulting empirical signed actor field has the score-function form
\citep{williams1992simple,sutton2000policy}:
\begin{equation}
    \mathbf{F}(\theta)
    =
    \mathbb{E}_{(s,a)\sim\nu}
    \left[
        \widehat{A}(s,a)
        \nabla_\theta \log \pi_\theta(a \mid s)
    \right],
    \label{eq:empirical_actor_field}
\end{equation}
with the corresponding update
\(\theta_{t+1}=\theta_t+\eta\mathbf{F}(\theta_t)\).
Repeated reuse means that the same historical samples and signed feedback
are evaluated under an evolving current policy;
Section~\ref{sec:theory} studies the actor-side dynamics induced by this
reuse.

To separate favorable and unfavorable feedback, write
\(\widehat A^{+}=\max\{\widehat A,0\}\) and
\(\widehat A^{-}=\max\{-\widehat A,0\}\), so the actor field decomposes as
\begin{equation}
\begin{aligned}
    \mathbf{F}(\theta)
    ={}&
    \mathbb{E}_{\nu}
    \left[
        \widehat{A}^{+}(s,a)
        \nabla_\theta \log \pi_\theta(a \mid s)
    \right] \\
    &-
    \mathbb{E}_{\nu}
    \left[
        \widehat{A}^{-}(s,a)
        \nabla_\theta \log \pi_\theta(a \mid s)
    \right].
\end{aligned}
\label{eq:signed_actor_decomposition}
\end{equation}
For a negative sample, the update magnitude factorizes as
$\widehat{A}^{-}(s,a)
\lVert\nabla_\theta \log \pi_\theta(a \mid s)\rVert$,
separating the severity of its estimated disadvantage from its
learner-relative policy geometry.

We use \emph{coefficient magnitude} for \(\widehat A^-\), \emph{score
response} for \(\|\nabla\log\pi\|\), \emph{update magnitude} for their
product, \emph{reuse loop gain} for the squared remoteness gradient
(Section~\ref{sec:distance_strength}), and \emph{aggregate force} for the
expectation or sum of signed updates.

All aggregate quantities use the same update base measure \(\nu\). We write
\begin{equation}
    p=\mathbb E_\nu[\widehat A^+],
    \qquad
    q=\mathbb E_\nu[\widehat A^-],
    \qquad
    \rho=q/p \quad (p>0).
    \label{eq:aggregate_masses}
\end{equation}
Because \(p\) and \(q\) are population expectations under \(\nu\), positive
and negative sample frequencies are already included.

% ------------------------------------------------------------------
% SECTION 3: THEORETICAL FRAMEWORK
% ------------------------------------------------------------------
\section{Far-Field Dynamics of Historical Policy Updates}
\label{sec:theory}

We first characterize how remoteness controls score contribution in
policy-output coordinates. We then show how reuse creates positive
self-attenuation and negative self-amplification, and derive the resulting
aggregate regimes and policy-family manifestations.

\subsection{Learner-Relative Remoteness and Update Strength}
\label{sec:distance_strength}

To isolate the action-side geometry, we fix a context $s$ and write
$u$ for the policy-output coordinates at that context. For a historical
action $a$, we define its learner-relative remoteness as
\begin{equation}
    D_u(a)
    =
    -\log \pi_u(a).
    \label{eq:learner_relative_remoteness}
\end{equation}
For a categorical policy, $D_u(a)$ is the surprisal of the selected
action. For a Gaussian policy, it equals a scale-dependent constant plus
one half of the squared Mahalanobis distance from the policy mean.
Thus, $D_u(a)$ measures how unlikely the historical action is under the
current learner, rather than defining a metric in the action space.

We measure the corresponding squared score norm by
\begin{equation}
    R_u(a)
    =
    \left\|
        \nabla_u \log \pi_u(a)
    \right\|_2^2.
    \label{eq:score_strength}
\end{equation}
A sample with advantage estimate $\widehat{A}(s,a)$ contributes an
output-space update of magnitude
$|\widehat{A}(s,a)|\sqrt{R_u(a)}$.
The following results characterize how learner-relative remoteness
controls this update strength.

\begin{proposition}[Distance-dependent score-function magnitude]
\label{prop:distance_strength}
Fix a context and a historical action \(a\).

For a Gaussian policy with mean \(\mu\) and fixed covariance
\(\Sigma\succ0\), let
\(C_\Sigma=\frac{d}{2}\log(2\pi)+\frac12\log|\Sigma|\) denote the
Gaussian normalization constant, which is independent of \(a\) and
\(\mu\). Then learner-relative remoteness satisfies
\begin{equation}
    D_\mu(a)-C_\Sigma
    =
    \frac12
    (a-\mu)^\top\Sigma^{-1}(a-\mu).
    \label{eq:gaussian_remoteness_mahalanobis}
\end{equation}
Moreover, the squared mean-score norm
\(R_\mu(a)=\|\nabla_\mu\log\pi_\mu(a)\|_2^2\) is bounded above and below
by positive constants times \(D_\mu(a)-C_\Sigma\), with constants
determined only by \(\Sigma\). Consequently, as standardized distance
increases, the Gaussian mean-score response grows without bound. When
\(\Sigma=\sigma^2I\), \(R_\mu(a)=2(D_\mu(a)-C_\Sigma)/\sigma^2\), so
remoteness and squared mean-score norm have the same global ordering.

For a categorical policy with logits \(z\), learner-relative remoteness is
the selected-action surprisal \(D_z(a)=-\log\pi_z(a)\). The selected-logit
score response is
\begin{equation}
    \left|
    \frac{\partial}{\partial z_a}
    \log \pi_z(a)
    \right|
    =
    1-e^{-D_z(a)}.
    \label{eq:categorical_selected_score}
\end{equation}
Thus, categorical surprisal can diverge as the selected probability
approaches zero, but the selected-logit response saturates at one. The
full logit-score norm remains bounded as well.
\end{proposition}

(Proof: Appendix~\ref{app:distance_strength}.) Gaussian remoteness can
produce an unbounded mean-score response, whereas categorical surprisal can
diverge while the logit-score response remains bounded.


\subsection{Self-Amplification under Historical Reuse}
\label{sec:self_amplification}

The preceding result compares samples at a fixed policy. We now consider
a fixed historical action that is repeatedly reused while the learner
changes. Write
\[
    D(u)=D_u(a),
    \qquad
    R(u)=\|\nabla_uD(u)\|_2^2,
\]
and consider the repeated single-sample update
\begin{equation}
    u_{t+1}
    =
    u_t
    +
    \eta\widehat A\nabla_u\log\pi_{u_t}(a)
    =
    u_t
    -
    \eta\widehat A\nabla_uD(u_t).
    \label{eq:single_sample_reuse_update}
\end{equation}
Here, the action \(a\) and its signed coefficient \(\widehat A\) are
held fixed over the reuse window.

\begin{theorem}[Self-attenuation and self-amplification under reuse]
\label{thm:reuse_dynamics}
Assume that \(D(u)\) is differentiable and convex on a convex
neighborhood containing the update segments, and define
\[
    D_t=D(u_t),
    \qquad
    R_t=R(u_t).
\]
If \(\widehat A<0\), then every reuse step satisfies
\begin{equation}
    D_{t+1}
    \ge
    D_t+\eta|\widehat A|R_t,
    \qquad
    R_{t+1}\ge R_t.
    \label{eq:negative_reuse_amplification}
\end{equation}
Thus, negative reuse makes the historical action increasingly remote
while preventing its score-function norm from decreasing.

If \(\widehat A>0\), and \(D\) additionally has an \(L\)-Lipschitz
gradient on the same neighborhood with \(0<\eta\widehat A\le 1/L\), then
\begin{equation}
    D_{t+1}
    \le
    D_t-\frac{\eta\widehat A}{2}R_t,
    \qquad
    R_{t+1}\le R_t.
    \label{eq:positive_reuse_attenuation}
\end{equation}
Thus, positive reuse makes the historical action increasingly compatible
with the learner while attenuating its subsequent score response.
\end{theorem}

Both policy coordinates satisfy the theorem's convexity assumption
globally: fixed-covariance Gaussian remoteness is quadratic in the mean,
and categorical negative log-likelihood is convex in logits
(Appendix~\ref{app:reuse_dynamics}).

The theorem separates a one-time distant update from a historical-reuse
feedback loop. Combined with Proposition~\ref{prop:distance_strength}, it
yields unbounded amplification for fixed-covariance Gaussian mean updates
and bounded but persistent amplification for categorical logits.

For an isolated isotropic-Gaussian negative sample, repeated
uncontrolled reuse produces geometric growth in remoteness; the exact
rate is given in Appendix~\ref{app:gaussian_reuse_rate}.

\subsection{Aggregate Equilibria under Historical Reuse}
\label{sec:aggregate_equilibria}

The preceding single-sample reuse analysis tracks one fixed historical
action. We now aggregate positive and negative historical samples and
ask whether their competing attraction and repulsion produce a finite
stable equilibrium, persistent drift, or instability.

To analyze continuous-action Gaussian and discrete-action categorical
policies within a single argument, we use a common exponential-family
output representation at a fixed context; the exact coordinates and
regularity assumptions are given in
Appendix~\ref{app:aggregate_equilibrium}. Using the aggregate masses
\(p,q,\rho\) from Equation~\eqref{eq:aggregate_masses}, let
\(\mathbf m_+\) denote the positive-only target, the mean-parameter point
selected by positive feedback alone. When \(q>0\), let \(\mathbf m_-\)
denote the magnitude-weighted negative-sample moment (a repulsion moment,
not an attraction target).

\begin{theorem}[Aggregate attraction--repulsion regimes]
\label{thm:aggregate_equilibria}
Consider the continuous aggregate actor dynamics and the corresponding
discrete update under the regular minimal exponential-family
representation specified in
Appendix~\ref{app:aggregate_equilibrium}. Assume \(p>0\).

\textbf{Positive-only case \((\rho=0)\).}
If \(\mathbf m_+\) lies in the interior of the attainable mean-parameter
space, there is a unique finite equilibrium whose mean parameter is
\(\mathbf m_+\). It is locally asymptotically stable in continuous time
and for sufficiently small discrete step sizes.

\textbf{Positive-dominant regime \((0<\rho<1)\).}
Define the signed target
\begin{equation}
    \mathbf m^\star(\rho)
    =
    \mathbf m_+
    +
    \frac{\rho}{1-\rho}
    \left(
        \mathbf m_+-\mathbf m_-
    \right).
    \label{eq:signed_target_ratio}
\end{equation}
If \(\mathbf m^\star(\rho)\) lies in the interior of the attainable
mean-parameter space, there is a unique finite equilibrium whose mean
parameter is \(\mathbf m^\star(\rho)\). It is locally asymptotically
stable in continuous time and for sufficiently small discrete step
sizes. Holding \(\mathbf m_+\) and \(\mathbf m_-\) fixed, its displacement
from \(\mathbf m_+\) grows without bound as \(\rho\to1^-\), unless
\(\mathbf m_+=\mathbf m_-\). If the signed target leaves the attainable
mean-parameter space, no finite equilibrium exists.

\textbf{Critical regime \((\rho=1)\).}
If \(\mathbf m_+\neq\mathbf m_-\), the policy-output coordinate exhibits
persistent drift and admits no finite equilibrium. If
\(\mathbf m_+=\mathbf m_-\), every point is stationary, but no point is
an isolated asymptotically stable equilibrium.

\textit{Negative-dominant regime (\(\rho>1\)).}
Every finite stationary point, if one exists, is unstable in continuous
time and under every discrete update with positive step size. Hence, this
regime admits no locally asymptotically stable finite equilibrium.
\end{theorem}

The proof, including the exponential-family actor field and the exact
discrete-time step-size condition, is provided in
Appendix~\ref{app:aggregate_equilibrium}.
Equation~\eqref{eq:signed_target_ratio} gives the central geometric
interpretation: while positive feedback remains dominant, negative
feedback displaces the equilibrium from \(\mathbf m_+\) in the direction
\(\mathbf m_+-\mathbf m_-\), and therefore away from
\(\mathbf m_-\).

% ------------------------------------------------------------------
% SECTION 4: METHODOLOGY
% ------------------------------------------------------------------
\subsection{Policy-Family Manifestations of Negative Dominance}
\label{sec:policy_family_manifestations}

Negative dominance admits no finite stable equilibrium, but its score
and support manifestations depend on the policy family. We therefore
analyze Gaussian and categorical policies separately.

\begin{theorem}[Policy-family score behavior under negative dominance]
\label{thm:policy_family_manifestations}
Assume \(p<q\).

\textbf{Gaussian policy.}
Consider a Gaussian policy
\(\pi_\mu=\mathcal N(\mu,\Sigma)\) with fixed
\(\Sigma\succ0\). Its aggregate mean dynamics have a unique finite
stationary point, which is repelling. Unless initialized exactly at that
point, the distance from it diverges under both the continuous dynamics
and every discrete trajectory with fixed step size \(\alpha>0\).
Consequently, for every fixed historical action \(a\),
\begin{equation}
    \left\|
        \nabla_\mu\log\pi_\mu(a)
    \right\|_2
    =
    \left\|
        \Sigma^{-1}(a-\mu)
    \right\|_2
    \longrightarrow
    \infty.
    \label{eq:gaussian_unbounded_score}
\end{equation}

\textbf{Categorical policy.}
Consider a categorical policy
\(\pi_z=\operatorname{softmax}(z)\) with at least two actions and the
zero-sum gauge \(\mathbf 1^\top z=0\). Any finite stationary point, when
one exists, is unique and repelling. Unless initialized exactly at that
point, the continuous dynamics and every discrete trajectory with fixed
step size \(\alpha>0\) eventually leave every compact subset of the
logit space. Along some subsequence, the corresponding probability
vector approaches the boundary of the simplex. Nevertheless, for every
action \(a\) and every finite \(z\),
\begin{equation}
\|\nabla_z\log\pi_z(a)\|_2^2
=
\|e_a-\pi_z\|_2^2
\le 2(1-\pi_z(a))^2
\le 2.
\label{eq:categorical_bounded_score}
\end{equation}
\end{theorem}

The proof is provided in
Appendix~\ref{app:policy_family_manifestations}.
Fixed covariance is sufficient for the Gaussian result; learned covariance is
not required.

\subsection{Why Selective Tapering Works}
\label{sec:selective_tapering_bridge}

The preceding regimes identify the failure; the following radial limit
identifies the control property that reverses it.

\begin{lemma}[Far-field radial balance]
\label{lem:far_field_radial_balance}
For a fixed-covariance Gaussian field with finite negative atoms,
\[
\begin{aligned}
F(\mu)&=p\Sigma^{-1}(m_+-\mu)
 +\sum_iq_i\Sigma^{-1}(\mu-a_i),\\
q&=\sum_iq_i,
\end{aligned}
\]
and \(\mu=m_++rv\), \(\|v\|_2=1\),
\[
\lim_{r\to\infty}\frac{\langle v,F(\mu)\rangle}{r}
=(q-p)v^\top\Sigma^{-1}v.
\]
\end{lemma}

\begin{corollary}[Selective tapering versus global scaling]
\label{cor:selective_vs_global}
If each negative atom is multiplied by
\(\omega_i(D_i)\to0\), the limit in
Lemma~\ref{lem:far_field_radial_balance} becomes
\(-p v^\top\Sigma^{-1}v<0\). Thus, for every fixed finite
\(\rho=q/p\), without \(\rho\)-specific tuning, the far-field radial
component is eventually inward; the onset radius is not uniform in
\(\rho\) and depends on the tail rate of \(\omega\). If additionally
\(\omega(D)\sqrt D\to0\), the absolute tapered negative force vanishes.
In contrast, global scaling \(\omega\equiv\alpha\) is inward only when
\(\alpha q<p\). Proofs are in Appendix~\ref{app:far_field_bridge}.
\end{corollary}

\begin{proposition}[Restoration by reference regularization]
\label{prop:regularized_restoration}
The regularized linear field has a unique finite stable equilibrium when
\(p>q\) for every \(\gamma\ge0\), when \(p=q\) for every \(\gamma>0\),
and when \(p<q\) if and only if \(\gamma>q-p\). At \(p=q\), its
displacement is \(O(\gamma^{-1})\). Reference regularization restores
stability by tethering the policy; DRPO instead removes the far-field
source. The proof and the Euclidean-spring variant are in
Appendix~\ref{app:regularized_restoration}.
\end{proposition}

Corollary~\ref{cor:selective_vs_global} justifies selective tapering: it
removes the remote source, whereas a constant global scale must be tuned to
the negative-to-positive mass ratio. These conclusions assume frozen data,
fixed advantages, and unbounded reuse; data refresh and actor--critic
coupling alter the field.

\section{Dynamic Remoteness-Aware Policy Optimization}
\label{sec:drpo}

Section~\ref{sec:selective_tapering_bridge} identifies the design target:
control the remote negative source without uniformly weakening the negative
branch. We first define DRPO, then derive its control-order hierarchy and
closed-loop guarantees.

\subsection{Update and Policy-Family Instantiations}
\label{subsec:drpo_update}

For remoteness \(D\), DRPO uses the thresholded exponential taper
\begin{equation}
    \omega_{\mathrm{Exp}}(D)
    =
    \exp
    \left\{
        -\lambda
        \left[
            \frac{D-\tau}{c}
        \right]_{+}
    \right\},
    \label{eq:drpo_exp_weight}
\end{equation}
where \(\tau\) marks the control threshold, \(c>0\) fixes the remoteness
scale, and \(\lambda>0\) controls the decay rate. The weight is exactly one
when \(D\leq\tau\) and decays smoothly when \(D>\tau\).

Using \(D_\theta(s,a)=-\log\pi_\theta(a\mid s)\), define
\(\bar D_\theta(s,a)=\operatorname{sg}[D_\theta(s,a)]\), where
\(\operatorname{sg}[\cdot]\) denotes stop-gradient. DRPO changes only the
negative branch:
\begin{equation}
    \widetilde A_{\theta}(s,a)
    =
    \widehat A^{+}(s,a)
    -
    \omega_{\mathrm{Exp}}
    \!\left(\bar D_\theta(s,a)\right)
    \widehat A^{-}(s,a),
    \label{eq:drpo_effective_advantage}
\end{equation}
with actor field
\begin{equation}
    \mathbf F_{\mathrm{DRPO}}(\theta)
    =
    \mathbb E_{(s,a)\sim\nu}
    \left[
        \widetilde A_{\theta}(s,a)
        \nabla_\theta\log\pi_\theta(a\mid s)
    \right].
    \label{eq:drpo_actor_field}
\end{equation}
Remoteness is recomputed under the current policy but detached before it is
used as a coefficient; DRPO is therefore a dynamic weighting rule rather than
a differentiable distance regularizer.

At each policy iterate, tapering preserves the aggregate field form
after replacing the negative mass and moment by
\[
q_\omega=\mathbb E_\nu[\widehat A^-\omega],
\qquad
\mathbf M_{-,\omega}=\mathbb E_\nu[\widehat A^-\omega T(a)].
\]
For \(0\le q_\omega<p\), the controlled signed target is
\[
\mathbf m_\omega^\star
=
\frac{p\mathbf m_+-\mathbf M_{-,\omega}}{p-q_\omega},
\]
with \(\mathbf M_{-,\omega}=0\) when \(q_\omega=0\).

\paragraph{Gaussian policies.}
For a Gaussian policy with mean \(\mu_\theta(s)\) and fixed covariance
\(\Sigma\), define
\begin{equation}
    D_{\mathrm G,\theta}(s,a)
    =
    (a-\mu_\theta(s))^\top
    \Sigma^{-1}
    (a-\mu_\theta(s)).
    \label{eq:gaussian_squared_distance}
\end{equation}
Since \(D_\theta-C_\Sigma=\frac12D_{\mathrm G,\theta}\), the additive
constant and factor \(1/2\) can be absorbed into the threshold and scale:
\begin{equation}
    \omega_{\mathrm G}(s,a)
    =
    \exp
    \left\{
        -\lambda_{\mathrm G}
        \left[
            \frac{D_{\mathrm G,\theta}(s,a)-\tau_{\mathrm G}}
                 {c_{\mathrm G}}
        \right]_{+}
    \right\}.
    \label{eq:drpo_gaussian_weight}
\end{equation}

\paragraph{Categorical policies.}
For a categorical policy, substituting selected-action surprisal into
Equation~\eqref{eq:drpo_exp_weight} gives
\begin{equation}
    \omega_{\mathrm C}(s,a)
    =
    \min
    \left\{
        1,\,
        \left(e^{\tau}\pi_\theta(a\mid s)\right)^{\lambda/c}
    \right\}.
    \label{eq:drpo_categorical_weight}
\end{equation}
\paragraph{Where existing off-policy controls intervene.}
Existing controls intervene at different points without sharing a design
goal. Importance weighting changes the influence of historical samples,
PPO-style methods constrain movement relative to a refreshed anchor
\citep{schulman2017proximal}, and reference regularization tethers the
policy. DRPO shares the sample-influence intervention point with importance
weighting, but replaces behavior-determined attenuation with a
learner-relative control law: remoteness is recomputed under the current
policy, the threshold and decay rate are designed, and only the negative
branch is weighted. With likelihood-defined remoteness, the clipped behavior
ratio \(\min\{(\pi_\theta/\mu)^\beta,1\}\) is itself an exponential taper,
with canonical TOPR as the \(\beta=1\) behavior-threshold case
\citep{leroux2025tapered}. Compared with movement- and location-based
constraints, DRPO directly suppresses the remote negative source while
preserving positive and near-field feedback. Exact identities,
anchor-refresh effects, and scope conditions are given in
Appendix~\ref{app:topr_relation}.

\subsection{Loop-Gain Order}
\label{subsec:min_order_control}

The quantity that drives the next change in remoteness is not the
single-step parameter displacement but the squared score
\(R(D)=\|\nabla D\|_2^2\). We call this the reuse \emph{loop gain}.

\begin{proposition}[Gaussian loop-gain order]
\label{prop:loop_gain_order}
For a fixed-covariance Gaussian mean policy,
\(I_{\mathrm{loop}}(D):=R(D)=\Theta(D)\) in the far field. Hence
\(\omega(D)I_{\mathrm{loop}}(D)\) is bounded exactly at the matching
order \(\omega(D)=O(D^{-1})\), and vanishes when
\(\omega(D)=o(D^{-1})\). Equivalently, the matching order is
\(O(r^{-2})\) in standardized distance \(r\), since
\(D-C_\Sigma=r^2/2\). The proof is in
Appendix~\ref{app:min_order_control}.
\end{proposition}

At the exact matching rate \(\omega(D)=\Theta(D^{-1})\), the weighted
parameter force is \(\Theta(D^{-1/2})\to0\), but the weighted reuse loop
gain is \(\Theta(1)\). An isolated reused atom can therefore keep driving
remoteness outward at a constant order. For isolated Gaussian reuse,
this yields geometric growth without control, linear growth under
matching-order reciprocal tapering, and logarithmic growth under
DRPO's exponential taper
(Appendix~\ref{app:gaussian_reuse_rate}).

\paragraph{Utility boundary.}
\label{subsec:distance_dependent_utility}
Remoteness is a control coordinate rather than a utility label. When
historical negative feedback no longer supplies local correction, DRPO
targets \(\omega(D)I(D)\to0\) as \(D\to\infty\); task direction remains
supplied by the underlying advantage or verifier signal. The corresponding
utility formulation is given in
Appendix~\ref{app:negative_update_utility}.

\subsection{Guarantee Hierarchy}
\label{sec:exponential_taper}

The DRPO taper satisfies the stronger requirement above for every
finite-order response considered here.

\begin{proposition}[Vanishing finite-order far-field response]
\label{prop:exp_vanishing_influence}
Let \(I(D)\) denote the raw far-field response scale. If, for some
finite \(m\geq0\) and \(C>0\),
\(I(D)\leq C(1+D)^m\) in the far field, then the taper in
Equation~\eqref{eq:drpo_exp_weight} satisfies
\[
    \omega_{\mathrm{Exp}}(D) I(D)
    \longrightarrow 0
    \qquad
    \text{as }D\to\infty .
\]
Thus, the exponential taper eliminates any finite-order far-field tail
rather than merely bounding it.
\end{proposition}

The exponential is a convenient, non-unique envelope; the proof and the
detailed L0--L3 guarantee hierarchy, including the mapping to global and
reciprocal controls, are given in
Appendix~\ref{app:exp_vanishing_influence}. The hierarchy ranks guarantees,
not task performance.

Closed-loop consequences are consistent with the same control principle. For
the fixed-covariance Gaussian mean field, exponential tapering makes trajectories
ultimately bounded and guarantees at least one finite equilibrium; if a stable
equilibrium lies entirely inside the untapered region, its local linearization
is unchanged. For a repeatedly suppressed categorical action, tapering changes
far-tail probability decay from exponential to polynomial rather than imposing
a nonzero probability floor. Full theorem statements, branch conditions, and
proofs are given in Appendices~\ref{app:gaussian_drpo_boundedness},
\ref{app:categorical_self_limiting}, and
\ref{app:categorical_restoring_balance}.

\section{Experiments}
\label{sec:experiments}

The experiments proceed from occurrence to identification, control, and
transfer. We
first test whether the predicted far-field pattern appears in realistic
continuous-control and autoregressive policies. Controlled environments then
separate coefficient magnitude from policy geometry, trace the resulting
update into policy behavior, and compare negative-weighting rules under known
near/far structure. The final task-level sections evaluate D4RL and
structured-generation tasks.

\subsection{Environments and Datasets}
\label{sec:environments_datasets}

We use C-U1 (a Gaussian continuous-action contextual bandit) and D-U1
(unordered categorical actions with shared representations) as controlled
environments with known ground truth for source isolation, policy dynamics,
and held-out-context generalization. External diagnostics use nine D4RL
locomotion cells~\citep{fu2020d4rl} and nine structured-generation tasks---
Countdown~\citep{hutton2002countdown,pan2025tinyzero}, seven Reasoning Gym
tasks~\citep{stojanovski2025reasoninggym}, and
WikiSQL~\citep{zhong2017seq2sql}---covering continuous-action and discrete
autoregressive policies; full configurations are in
Appendix~\ref{app:experimental_details}.

\paragraph{Conventions.}
Controlled source-isolation reports policy-score response, whereas neural
probes report implemented actor-gradient magnitude to test whether the
near-to-far pattern persists in trainable updates. Task-performance collapse,
support/variance-boundary events, and NaN/Inf failures remain separate
outcomes (Appendix~\ref{app:controlled_protocols}).

\subsection{External Diagnostic of Far-Field Repulsion}
\label{sec:external_far_field}

For D4RL and structured-generation tasks, learner-relative remoteness varies
while sample-quality magnitude is held fixed; full protocols are given in
Appendices~\ref{app:d4rl_gradient_diagnostic}
and~\ref{app:structured_generation_methods}.

Across both D4RL and structured-generation tasks, the implemented
actor-gradient increases from the near to the far field
(Figure~\ref{fig:app_external_gradient_panels}). The approximately linear
D4RL trend agrees with the Gaussian prediction, while the fixed-coefficient
structured-generation trend is consistent with persistent categorical
suppression. These probes show that the predicted near-to-far pattern appears
in realistic neural policies; the controlled experiments below separate
remoteness from coefficient magnitude and test whether far-field updates
drive instability.




\subsection{Controlled Identification of the Causal Mechanism}
\label{sec:controlled_identification}

Because external diagnostics entangle sample quality, remoteness,
critic/verifier coefficients, coverage, and downstream dynamics, the controlled
experiments vary coefficient magnitude and remoteness independently and
intervene on near- versus far-field negative updates
(Appendix~\ref{app:controlled_protocols}). In C-U1, negative actions share a
reward contour and fixed advantage but differ in standardized distance; the
categorical direct-softmax reference fixes the negative coefficient while
varying current-policy surprisal. For the causal intervention, all C-U1
branches share data, fixed labels, initialization, optimizer, and horizon and
differ only in near/far negative-update treatment
(Appendix~\ref{app:controlled_shared}).

\begin{figure}[t]
    \centering
    \includegraphics[
        width=\columnwidth
    ]{figures/figure1_external_gradient_d4rl9.pdf}

    \caption{
\textbf{External far-field diagnostics.}
Relative implemented actor-gradient versus learner-relative remoteness,
aggregated over nine D4RL locomotion dataset cells in (a) and
structured-generation tasks in (b). Gradients are normalized to the
corresponding near-field reference; dashed lines show matched negative-signal
controls and shading shows 95\% confidence intervals
(disaggregated panels: Figure~\ref{fig:app_external_gradient_panels}).
}
    \Description{Two-panel external diagnostic of relative implemented
actor-gradient versus learner-relative remoteness in D4RL and
structured-generation tasks.}
    \label{fig:external_far_field}
\end{figure}

\subsubsection{Remoteness as an Independent Source of Update Magnitude}
\label{sec:controlled_source}

\begin{figure*}[t]
\centering

\begin{minipage}[t]{0.49\textwidth}
    \centering
    \vspace{0pt}

    \includegraphics[
        width=0.92\linewidth,
        trim={0pt 2pt 0pt 2pt},
        clip
    ]{figures/fig_6_3_1_source_heatmap.pdf}

    \vspace{0.35em}

    \includegraphics[
        width=0.92\linewidth,
        trim={0pt 2pt 0pt 2pt},
        clip
    ]{figures/fig_6_3_2_rescue_plot.pdf}

    \vspace{-0.4em}

\caption{
    \textbf{Controlled source isolation and causal transmission.}
    \textbf{Top:} coefficient magnitude (COE) and policy-score response (SC)
    across remoteness bins, normalized to the near bin, for the C-U1 Gaussian
    protocol and the direct-softmax categorical reference.
    \textbf{Bottom:} C-U1 interventions; Ret. denotes reward retention
    (20-seed means with bootstrap CIs), and Task coll. denotes task-performance
    collapse counts.
}
    \label{fig:controlled_6_3}
\end{minipage}
\hfill
\begin{minipage}[t]{0.49\textwidth}
    \centering
    \vspace{0pt}

    \includegraphics[
        width=0.86\linewidth
    ]{figures/fig_6_4_1_phase_transition.pdf}

    \vspace{-0.35em}
\caption{
    \textbf{Stable extrapolation from controlled negative feedback.}
    Increasing \(q/p\) first lifts held-out-context reward above the
    positive-only ceiling, then causes over-extrapolation and task-performance
    collapse; policy shift measures displacement from the positive-only target.
    The strip reports task collapse, boundary events, and NaN/Inf failures
    separately.
}
    \label{fig:phase_transition_6_4_1}
\end{minipage}

\end{figure*}

The top panel of Figure~\ref{fig:controlled_6_3} separates the
coefficient term from the policy-score term. The coefficient magnitude
remains flat across remoteness bins, whereas the score response increases
with remoteness. Thus, the larger update cannot be explained
by farther samples having larger advantages; it is generated by the
policy-score geometry itself. The manifestation differs across action
spaces: Gaussian policies exhibit amplitude amplification, while
categorical policies approach a bounded direct-logit score but continue
to suppress the selected action toward the simplex boundary. The source-isolation experiment establishes where the update difference comes
from; the next intervention tests whether it propagates into drift and
collapse.

\subsubsection{Causal Transmission into Drift and Collapse}
\label{sec:controlled_transmission}

The bottom panel of Figure~\ref{fig:controlled_6_3} shows a clear rescue
pattern. Removing near-field negative updates leaves performance collapse
largely unchanged, whereas removing or capping far-field updates restores
reward retention and eliminates task-performance collapse. Global
rescaling also rescues the run, supporting the interpretation that
abnormal negative-update magnitude is a mediator. In contrast,
transferring the far-field budget to near-field negatives gives only
partial recovery, showing that the rescue is not obtained by simply
preserving the same total negative-update budget.

These interventions identify the far-field negative update, rather than
negative feedback in general, as a causal transmission path to policy
instability in the controlled Gaussian environment.

\subsection{From Useful Repulsion to DRPO Control}
\label{sec:repulsion_control}

The causal rescue does not imply that all negative feedback should be removed.
C-U1 and D-U1 place a hidden optimum beyond the positive-only target and make
local repulsion point toward it, while additional remote negatives create
far-field pressure. This construction separates useful repulsion from harmful
dominance and permits controlled comparisons of taper shape.

\subsubsection{Stable Extrapolation from Controlled Negative Feedback}
\label{sec:controlled_negative_strength}



We sweep the effective negative strength relative to positive attraction,
denoted by $q/p$, under the matched-control invariants above; protocol and
metric details are in Appendix~\ref{app:negative_strength_sweep}. The
positive-only setting is $q/p=0$.

Figure~\ref{fig:phase_transition_6_4_1} shows a non-monotone transition.
For small and moderate $q/p$, held-out-context reward rises above the
positive-only ceiling, showing that controlled negative feedback provides
useful repulsion beyond the positive-support solution. As $q/p$ grows
larger, policy shift continues to increase but reward declines,
indicating over-extrapolation rather than additional useful
generalization. At the largest strengths, task-performance collapse and
boundary events appear; numerical failures remain a separate outcome class.

Positive-only learning removes the harmful tail but also the beneficial
intermediate regime; remoteness-aware control should preserve local repulsion
while suppressing the far-field tail.


Matched controls confirm the intended selectivity: remoteness-aware tapering can
preserve near-field repulsion while suppressing far-field influence, whereas
global scaling reduces both indiscriminately. Detailed comparisons of
Positive-only, Global-\(\alpha\), Reciprocal-Linear, Reciprocal-Quadratic, and
exponential tapering under near-field-retention and \(B(\omega)\)-matched
controls are reported in Appendix~\ref{app:controlled_taper_protocol}.

\subsubsection{External Validation of Remoteness Tapering}
\label{sec:external_tapers}

\begin{wraptable}{r}{0.52\columnwidth}
\vspace{-0.7em}
\centering
\caption{\textbf{Protocol-matched external transfer.}
D4RL aggregate and Countdown-0.5B late-window Pass@8; Countdown reports the
best observed two-seed development point for each method.}
\label{tab:taper_transfer_summary}
\scriptsize
\renewcommand{\arraystretch}{0.88}
\resizebox{\linewidth}{!}{%
\begin{tabular}{lccccc}
\toprule
 & Pos. & G-$\alpha$ & Rec.-L & Rec.-Q & DRPO \\
\midrule
D4RL & 70.1 & 72.9 & 74.3 & 72.1 & \textbf{77.7} \\
Cnt. & 14.0 & 14.2 & 13.6 & 14.5 & \textbf{15.7} \\
\bottomrule
\end{tabular}%
}
\vspace{-0.5em}
\end{wraptable}

The same principle is tested externally: DRPO weights D4RL locomotion transitions by
standardized remoteness from the current Gaussian policy and Countdown
responses by mean-token surprisal, with protocol-matched methods sharing
data, labels, initialization, optimizer, and budget. Because external tasks
lack counterfactual near/far geometry, Table~\ref{tab:taper_transfer_summary}
provides transfer evidence rather than a closed cross-task ranking.

\subsection{Overall Task Performance}
\label{sec:overall_performance}

We evaluate task-level transfer on D4RL locomotion and structured generation,
reporting normalized return and verifier success, respectively.
Both evaluations train from fixed data, while learner-relative remoteness is
recomputed under the current policy during optimization. Protocol-matched
controls share data and compute budgets; task-performance collapse,
support/variance-boundary events, and NaN/Inf failures remain separate outcomes.

% Required packages:
% \usepackage{booktabs}
% \usepackage{graphicx}

\subsubsection{D4RL Locomotion}
\label{sec:d4rl_performance}

For continuous control, we evaluate DRPO on D4RL
\citep{fu2020d4rl}, a widely used offline-RL benchmark suite.
Table~\ref{tab:d4rl_performance} compares DRPO with published
CQL~\citep{kumar2020conservative}, TD3+BC~\citep{fujimoto2021minimalist},
and IQL~\citep{kostrikov2021offline} reference results; protocol and
implementation details are provided in Appendix~\ref{app:d4rl_methods}.

Without adding an explicit behavior-policy constraint or a conservative
value penalty on top of the shared advantage-weighted actor form, DRPO
remains competitive on the D4RL-9 total; its only method-specific
component is a detached remoteness taper on the negative branch. This is
consistent with the far-field source-control route of
Section~\ref{sec:selective_tapering_bridge}, but should be read as
task-level transfer rather than a causal test of the mechanism.

\subsubsection{Structured Generation}
\label{sec:structured_generation_performance}

For structured generation, we evaluate Qwen2.5-0.5B-Instruct
\citep{qwen2024qwen25} on the nine tasks in
Section~\ref{sec:environments_datasets}, using deterministic verifiers
and Pass@8. For each prompt, we construct a fixed bank of verified correct
and incorrect completions and reuse it throughout training, yielding a
zero-refresh stale-negative-reuse setting
(Appendices~\ref{app:structured_generation_dataset}
and~\ref{app:structured_generation_methods}).

We compare DRPO with AsymRE~\citep{arnal2025asymmetric}, joint
fitted-reference \(\beta\)-TOPR~\citep{leroux2025tapered}, and canonical
DPO~\citep{rafailov2023direct}; method-specific initialization is reported
in Table~\ref{tab:structured_generation_performance} and
Appendix~\ref{app:structured_generation_methods}.

\begin{table*}[t]
\centering
\caption{Normalized returns on the D4RL-9 Gym locomotion benchmark in (a)
and Pass@8 on structured-generation tasks in (b).}
\label{tab:external_task_performance}
\label{tab:d4rl_performance}
\label{tab:structured_generation_performance}
\vspace{-2pt}

\begin{minipage}[t]{0.465\textwidth}
\centering
\textbf{(a) D4RL-9 normalized return}\par\vspace{2pt}
{\scriptsize
\setlength{\tabcolsep}{2.05pt}
\renewcommand{\arraystretch}{1.04}
\begin{tabular}{@{}lrrrr@{}}
\toprule
Dataset & CQL$^\dagger$ & TD3+BC$^\dagger$ & IQL$^\dagger$ & \textbf{DRPO} \\
\midrule
halfcheetah-medium & 44.0 & \textbf{48.3} & \underline{47.4} & 47.3 \\
hopper-medium & 58.5 & \underline{59.3} & \textbf{66.3} & 57.3 \\
walker2d-medium & 72.5 & \underline{83.7} & 78.3 & \textbf{85.5} \\
halfcheetah-replay & \textbf{45.5} & \underline{44.6} & 44.2 & 43.3 \\
hopper-replay & \textbf{95.0} & 60.9 & \underline{94.7} & 87.8 \\
walker2d-replay & 77.2 & \textbf{81.8} & 73.9 & \underline{78.1} \\
halfcheetah-expert & \textbf{91.6} & 90.7 & 86.7 & \underline{91.1} \\
hopper-expert & \textbf{105.4} & \underline{98.0} & 91.5 & \underline{98.0} \\
walker2d-expert & 108.8 & \underline{110.1} & 109.6 & \textbf{110.4} \\
\midrule
\textit{D4RL-9 total} & \textit{\underline{698.5}} & \textit{677.4} &
\textit{692.4} & \textbf{\textit{698.8}} \\
\bottomrule
\end{tabular}
}
\end{minipage}
\hfill
\begin{minipage}[t]{0.515\textwidth}
\centering
\textbf{(b) Structured generation Pass@8}\par\vspace{2pt}
{\scriptsize
\setlength{\tabcolsep}{2.1pt}
\renewcommand{\arraystretch}{1.04}
\begin{tabular}{@{}lrrrrr@{}}
\toprule
Task & PO & AsymRE & $\beta$-TOPR & DPO$^\ddagger$ & \textbf{DRPO} \\
\midrule
Countdown & 13.98 & 14.38 & 14.41 & & \textbf{15.72} \\
Word Sorting & 62.06 & 61.72 & 59.61 & & \textbf{63.38} \\
Spiral Matrix & \textbf{100.00} & 99.92 & \textbf{100.00} & & 99.94 \\
Mini Sudoku & 95.91 & \textbf{97.58} & 96.88 & & 97.06 \\
Maze & 76.74 & \textbf{77.81} & 75.86 & & 77.66 \\
Word Ladder & 10.68 & 12.50 & 11.41 & & \textbf{13.38} \\
Knights \& Knaves & 67.53 & \textbf{70.16} & 68.59 & & 69.53 \\
Graph Coloring & 98.67 & 98.91 & 99.37 & & \textbf{99.88} \\
WikiSQL & 85.65 & 85.78 & 82.89 & & \textbf{86.69} \\
\midrule
\textit{Macro Avg.} & \textit{67.91} & \textit{68.75} & \textit{67.67} &
& \textbf{\textit{69.25}} \\
\bottomrule
\end{tabular}
}
\end{minipage}

\vspace{3pt}
\begin{minipage}{0.985\textwidth}
\scriptsize
\textit{Note.} PO = Positive-only.
\textbf{(a)} ``replay'' and ``expert'' abbreviate ``medium-replay'' and
``medium-expert'', respectively. $\dagger$: published results rather than
protocol-matched reruns. Bold/underline: best/second-best among the methods
shown. DRPO reports the mean over 10 runs.
\textbf{(b)} AsymRE, $\beta$-TOPR, and DRPO initialize LoRA directly from
Qwen2.5-0.5B-Instruct; $^{\ddagger}$DPO uses the same LoRA parameterization
but starts its policy and frozen reference from the short SFT checkpoint
(Appendix~\ref{app:structured_generation_methods}).
$\beta$-TOPR denotes the joint fitted-reference variant of TOPR
\citep{leroux2025tapered}.
For DRPO, the reported transfer-task scores are the maximum means over
contiguous five-point windows of the registered $\lambda$-response curves.
Macro Avg. is the unweighted mean of the nine task-level Pass@8 percentages.
\end{minipage}
\end{table*}

Across the eight transfer tasks, the reported DRPO score exceeds the
reported AsymRE score on five tasks and the Joint Fitted-Reference
\(\beta\)-TOPR score on seven tasks.

\section{Conclusion}

Repeated off-policy reuse can turn useful negative feedback into a
self-reinforcing far-field effect, linking per-sample repulsion to stable
displacement, persistent drift, and instability. DRPO selectively attenuates
remote negative updates while preserving local ones, and matched diagnostics
and interventions support this mechanism across continuous and discrete
policies. Our guarantees analyze the zero-refresh limit, which cleanly isolates
these dynamics, while the same learner-relative weighting principle also
applies when negative samples are already remote under the current policy,
including replay and stale-policy settings. Thus, remoteness-aware control is
not inherently restricted to static offline datasets.


% Acknowledgements should only appear in the accepted version.


\bibliography{example_paper,missing_references}
\bibliographystyle{icml2026}

%%%%%%%%%%%%%%%%%%%%%%%%%%%%%%%%%%%%%%%%%%%%%%%%%%%%%%%%%%%%%%%%%%%%%%%%%%%%%%%
%%%%%%%%%%%%%%%%%%%%%%%%%%%%%%%%%%%%%%%%%%%%%%%%%%%%%%%%%%%%%%%%%%%%%%%%%%%%%%%
% APPENDIX
%%%%%%%%%%%%%%%%%%%%%%%%%%%%%%%%%%%%%%%%%%%%%%%%%%%%%%%%%%%%%%%%%%%%%%%%%%%%%%%
%%%%%%%%%%%%%%%%%%%%%%%%%%%%%%%%%%%%%%%%%%%%%%%%%%%%%%%%%%%%%%%%%%%%%%%%%%%%%%%
\newpage
\appendix
\onecolumn
\section{Proofs for the Theoretical Results}
\label{app:theory_proofs}

\paragraph{Proof map.}
The score and reuse results are proved in
Sections~\ref{app:distance_strength}--\ref{app:gaussian_reuse_rate}; the
aggregate regimes and policy-family manifestations are proved in
Sections~\ref{app:aggregate_equilibrium} and
\ref{app:policy_family_manifestations}. The selective-tapering bridge,
reference restoration, Gaussian boundedness, control order, exponential tail,
and categorical results appear in Sections~\ref{app:far_field_bridge},
\ref{app:regularized_restoration}, \ref{app:gaussian_drpo_boundedness},
\ref{app:min_order_control}, \ref{app:exp_vanishing_influence}, and
\ref{app:categorical_self_limiting}--\ref{app:categorical_restoring_balance},
respectively.

\subsection{Proof of Proposition~\ref{prop:distance_strength}}
\label{app:distance_strength}

\begin{proof}
Let \(\delta=a-\mu\) and \(x=\Sigma^{-1/2}\delta\). For a fixed-covariance
Gaussian policy,
\[
D_\mu(a)-C_\Sigma
=
\frac12\|x\|_2^2,
\qquad
\nabla_\mu\log\pi_\mu(a)
=
\Sigma^{-1}\delta,
\qquad
R_\mu(a)
=
x^\top\Sigma^{-1}x.
\]
The Rayleigh quotient therefore gives
\[
2\lambda_{\min}(\Sigma^{-1})
\bigl(D_\mu(a)-C_\Sigma\bigr)
\le
R_\mu(a)
\le
2\lambda_{\max}(\Sigma^{-1})
\bigl(D_\mu(a)-C_\Sigma\bigr).
\]
Hence the Gaussian mean-score response is linearly comparable to
\(D_\mu(a)-C_\Sigma\) and grows without bound with standardized distance.
If \(\Sigma=\sigma^2I\), the bounds collapse to
\[
R_\mu(a)
=
\frac{2}{\sigma^2}
\bigl(D_\mu(a)-C_\Sigma\bigr).
\]

For a categorical policy, the \(K=1\) case is degenerate:
\(\pi_z(a)=1\), \(D_z(a)=0\), and the score is identically zero.
Assume \(K\ge2\) below and let
\(p_a=\pi_z(a)=e^{-D_z(a)}\). Since
\[
\nabla_z\log\pi_z(a)=e_a-\pi_z,
\]
the selected-logit response is
\[
\left|
\frac{\partial}{\partial z_a}\log\pi_z(a)
\right|
=
1-p_a
=
1-e^{-D_z(a)},
\]
which increases with surprisal and saturates at one. Moreover,
\[
R_z(a)
=
(1-p_a)^2+\sum_{j\ne a}\pi_z(j)^2.
\]
The non-selected probabilities are nonnegative and sum to \(1-p_a\), so
\[
\frac{(1-p_a)^2}{K-1}
\le
\sum_{j\ne a}\pi_z(j)^2
\le
(1-p_a)^2.
\]
Therefore
\[
\frac{K}{K-1}(1-p_a)^2
\le
R_z(a)
\le
2(1-p_a)^2
\le
2,
\]
which proves the claimed bounded categorical score behavior.
\end{proof}

\subsection{Proof of Theorem~\ref{thm:reuse_dynamics}}
\label{app:reuse_dynamics}

For fixed-covariance Gaussian means,
\[
D_\mu(a)
=
C_\Sigma+\frac12(a-\mu)^\top\Sigma^{-1}(a-\mu),
\qquad
\nabla_\mu^2D_\mu(a)=\Sigma^{-1}\succeq0.
\]
For categorical logits \(z\), with \(\pi=\operatorname{softmax}(z)\),
\[
D_z(a)
=
\log\sum_j e^{z_j}-z_a,
\qquad
\nabla_z^2D_z(a)
=
\operatorname{diag}(\pi)-\pi\pi^\top\succeq0,
\]
because
\(v^\top(\operatorname{diag}(\pi)-\pi\pi^\top)v
=\operatorname{Var}_{i\sim\pi}(v_i)\ge0\).
Thus both policy-output coordinates satisfy the theorem's convexity
assumption.

\begin{proof}
Write
\[
f=D,
\qquad
g_t=\nabla f(u_t),
\qquad
R_t=\|g_t\|_2^2.
\]

If \(\widehat A=-c<0\), let \(\alpha=\eta c\). Then
\(u_{t+1}=u_t+\alpha g_t\), and convexity gives
\[
D_{t+1}
\ge
D_t+g_t^\top(u_{t+1}-u_t)
=
D_t+\alpha R_t.
\]
Gradient monotonicity gives
\[
(g_{t+1}-g_t)^\top(u_{t+1}-u_t)\ge0.
\]
Since \(u_{t+1}-u_t=\alpha g_t\),
\(g_{t+1}^\top g_t\ge\|g_t\|_2^2\); Cauchy--Schwarz then yields
\(\|g_{t+1}\|_2\ge\|g_t\|_2\), with the zero-gradient case immediate.
Hence
\[
D_{t+1}
\ge
D_t+\eta|\widehat A|R_t,
\qquad
R_{t+1}\ge R_t.
\]
The second inequality is strict whenever the gradient changes
nontrivially along a direction of positive curvature.

If \(\widehat A=c>0\), again let \(\alpha=\eta c\), so
\(u_{t+1}=u_t-\alpha g_t\). The descent lemma gives
\[
D_{t+1}
\le
D_t-\alpha\left(1-\frac{\alpha L}{2}\right)R_t
\le
D_t-\frac{\alpha}{2}R_t
\]
whenever \(0<\alpha\le1/L\). For a convex \(L\)-smooth function,
cocoercivity gives
\[
(g_t-g_{t+1})^\top(u_t-u_{t+1})
\ge
\frac1L\|g_t-g_{t+1}\|_2^2.
\]
Let \(d=g_{t+1}-g_t\). Since \(u_t-u_{t+1}=\alpha g_t\),
\[
d^\top g_t
\le
-\frac{\|d\|_2^2}{\alpha L},
\]
and therefore
\[
R_{t+1}
=
\|g_t+d\|_2^2
\le
R_t+
\left(1-\frac{2}{\alpha L}\right)\|d\|_2^2
\le
R_t.
\]
Substituting \(\alpha=\eta\widehat A\) proves the positive-reuse claim.
\end{proof}

\subsection{Quantitative Gaussian Reuse Rates}
\label{app:gaussian_reuse_rate}

\begin{proposition}[Quantitative Gaussian reuse rates]
\label{prop:reuse_rate}
For an isotropic Gaussian mean policy with covariance \(\sigma^2I\),
one fixed negative sample of mass \(q>0\), and
\(X_t=D_t-C_\Sigma>0\), the uncontrolled update satisfies
\[
 X_{t+1}=\left(1+\eta q/\sigma^2\right)^2X_t,
\]
so remoteness grows geometrically. If the same sample is weighted by
\(\exp\{-a[D_t-\tau]_+\}\), \(a>0\), then its continuous-time
remoteness is \(O(a^{-1}\log t)\), and the discrete update obeys
\[
 X_t\le a^{-1}\{\log t+\log\log t+O(1)\}.
\]
This slowdown does not by itself prove aggregate boundedness; that requires
the closed-loop attraction established in
Theorem~\ref{thm:gaussian_drpo_boundedness}.
\end{proposition}

\begin{proof}[Proof of Proposition~\ref{prop:reuse_rate}]
Let \(\delta_t=\mu_t-a\) and
\(X_t=\|\delta_t\|_2^2/(2\sigma^2)=D_t-C_\Sigma\).
For an uncontrolled negative sample of mass \(q\),
\[
\delta_{t+1}
=\left(1+\frac{\eta q}{\sigma^2}\right)\delta_t,
\]
which gives the stated exact geometric recurrence.

Now let \(\omega_t=e^{-a[D_t-\tau]_+}\). Once
\(D_t>\tau\), write
\(b=\eta q/\sigma^2\) and
\(K_0=e^{a(\tau-C_\Sigma)}\). Then
\begin{equation}
X_{t+1}=X_t\left(1+bK_0e^{-aX_t}\right)^2.
\label{eq:app_tapered_gaussian_recurrence}
\end{equation}
Consequently \(\Delta X_t\ge0\), and for all sufficiently large
\(X_t\),
\begin{equation}
0\le\Delta X_t\le C X_te^{-aX_t}
\label{eq:app_gaussian_increment_bound}
\end{equation}
for a finite constant \(C\). Put \(Y_t=e^{aX_t}\). Since
\(aX_te^{-aX_t}\le1/e\),
Equation~\eqref{eq:app_gaussian_increment_bound} also gives a uniform
bound \(a\Delta X_t\le M\). Hence
\begin{align}
\Delta Y_t
&=Y_t(e^{a\Delta X_t}-1) \\
&\le C_1Y_t\,a\Delta X_t
\le C_2\log Y_t.
\label{eq:app_gaussian_y_increment}
\end{align}
The standard comparison sequence
\(\overline Y_t=K t\log t\), with \(K\) chosen larger than the
constant in Equation~\eqref{eq:app_gaussian_y_increment}, is eventually
a supersolution. Therefore
\(Y_t=O(t\log t)\) and
\[
X_t\le a^{-1}\{\log t+\log\log t+O(1)\}.
\]
If an application does not make the sign of \(\Delta X_t\) explicit,
negative increments already satisfy the desired upper comparison, while
the nonnegative branch obeys the argument above.

In continuous time, the far-branch equation is
\(\dot X=C_3Xe^{-aX}\). Thus
\(Y=e^{aX}\) satisfies \(\dot Y=C_3\log Y\), and the same comparison
gives \(X(t)=O(a^{-1}\log t)\).
\end{proof}

\subsection{Far-Field Radial Bridge}
\label{app:far_field_bridge}

\begin{proof}[Proofs of Lemma~\ref{lem:far_field_radial_balance} and
Corollary~\ref{cor:selective_vs_global}]
Set \(\mu=m_++rv\) with \(\|v\|_2=1\). Direct expansion gives
\begin{align}
\frac{\langle v,F(\mu)\rangle}{r}
={}&-p\,v^\top\Sigma^{-1}v
+\sum_iq_i v^\top\Sigma^{-1}v \\
&+\frac1r\sum_iq_i
\langle v,\Sigma^{-1}(m_+-a_i)\rangle.
\end{align}
The final term vanishes, proving the lemma.

With atom-specific weights, \(D_\mu(a_i)\to\infty\) along every ray.
For finitely many atoms, \(\omega_i(D_i)\to0\) therefore makes both the
weighted coefficient of \(r\) and the weighted constant term vanish.
The radial limit is
\(-p v^\top\Sigma^{-1}v\), which is strictly negative. This is a
pointwise-in-\(q/p\) statement: the radius at which it becomes negative
can grow with \(q/p\). Moreover,
\(\|\Sigma^{-1}(\mu-a_i)\|_2=O(\sqrt{D_i})\), so
\(\omega(D)\sqrt D\to0\) removes the absolute negative force. Finally,
for \(\omega_i\equiv\alpha\), the lemma applies with \(q\) replaced by
\(\alpha q\), yielding the condition \(\alpha q<p\).
\end{proof}

\subsection{Reference-Regularized Restoration}
\label{app:regularized_restoration}

\begin{proof}[Proof of Proposition~\ref{prop:regularized_restoration}]
The unweighted Gaussian field with a KL-geometry spring is
\[
F_\gamma(\mu)=\Sigma^{-1}
\{(q-p-\gamma)\mu+pm_+-qm_-+\gamma\mu_0\}.
\]
Whenever \(p-q+\gamma\ne0\), its unique equilibrium is
\[
\mu_\gamma^*
=\frac{pm_+-qm_-+\gamma\mu_0}{p-q+\gamma},
\]
and the field Jacobian is
\((q-p-\gamma)\Sigma^{-1}\). It is Hurwitz exactly when
\(\gamma>q-p\). This yields all three regimes in the proposition; for
\(p=q\), the displayed equilibrium has displacement proportional to
\(1/\gamma\). When \(p>q\), increasing \(\gamma\) contracts the
equilibrium toward \(\mu_0\), potentially sacrificing extrapolation.

For the Euclidean spring \(-\gamma(\mu-\mu_0)\), the Jacobian is
\((q-p)\Sigma^{-1}-\gamma I\). In the negative-dominant regime it is
Hurwitz exactly when
\(\gamma>(q-p)\lambda_{\max}(\Sigma^{-1})\); this condition is not
interchangeable with the KL-geometry coefficient above.
\end{proof}

\subsection{Gaussian Closed-Loop Boundedness}
\label{app:gaussian_drpo_boundedness}


\begin{theorem}[Gaussian DRPO boundedness]
\label{thm:gaussian_drpo_boundedness}
Let \(p>0\), \(\Sigma\succ0\), and
\[
F_\omega(\mu)=p\Sigma^{-1}(m_+-\mu)
+\sum_iq_i\omega_i(D_\mu(a_i))\Sigma^{-1}(\mu-a_i).
\]
For finitely many negative atoms with exponential tapers, define
\begin{equation}
\begin{aligned}
B_\omega&=\sup_\mu\left\|\sum_iq_i\omega_i(D_i)
\Sigma^{-1}(\mu-a_i)\right\|_2,\\
R^*&=\frac{B_\omega}{p\lambda_{\min}(\Sigma^{-1})}.
\end{aligned}
\label{eq:gaussian_bound_radius}
\end{equation}
Every continuous trajectory satisfies
\(\limsup_{t\to\infty}\|\mu(t)-m_+\|_2\le R^*\), and the field has at
least one finite equilibrium. If an uncontrolled stable equilibrium
places every negative atom strictly below its threshold, the controlled
and original fields have the same local linearization in a neighborhood of
that equilibrium.
\end{theorem}

The branch-wise bound below includes the case where the shifted threshold is
negative and tapering is active from radius zero. Along every escaping
sequence,
\[
\rho_{\mathrm{eff}}(\mu)
=
p^{-1}\sum_i q_i\omega_i(D_i)
\longrightarrow 0;
\]
this explains how the remote negative source disappears, but it is a
diagnostic consequence rather than the boundedness argument.

\begin{proof}[Proof of Theorem~\ref{thm:gaussian_drpo_boundedness}]
Let \(e=\mu-m_+\) and
\[
N_\omega(\mu)=\sum_iq_i\omega_i(D_i)
\Sigma^{-1}(\mu-a_i).
\]
The finite-atom exponential-taper field is locally Lipschitz and has at
most linear growth, so its solutions exist for all forward times.
For \(V(e)=\frac12\|e\|_2^2\),
\begin{align}
\dot V
&=-p e^\top\Sigma^{-1}e+e^\top N_\omega(\mu) \\
&\le-p\lambda_{\min}(\Sigma^{-1})\|e\|_2^2
+B_\omega\|e\|_2.
\label{eq:app_gaussian_lyapunov_bound}
\end{align}
Thus \(\dot V<0\) outside every ball of radius
\(R^*+\varepsilon\). The standard ultimate-boundedness argument gives
\[
\limsup_{t\to\infty}\|\mu(t)-m_+\|_2\le R^*.
\]
On the boundary of any ball centered at \(m_+\) with radius greater
than \(R^*\), Equation~\eqref{eq:app_gaussian_lyapunov_bound} is
strictly negative. Continuity of the field and the Poincar\'e--Bohl
inward-pointing theorem therefore imply at least one zero inside the
ball. This is an existence result, not a uniqueness or global-convergence
claim.

It remains to verify \(B_\omega<\infty\) and the stated constant. Put
\(r_i=\|\Sigma^{-1/2}(\mu-a_i)\|_2\), so
\(D_i=C_\Sigma+r_i^2/2\), and
\[
\|\Sigma^{-1}(\mu-a_i)\|_2
\le\sqrt{\lambda_{\max}(\Sigma^{-1})}\,r_i.
\]
For \(a=\lambda/c\) and
\(\widetilde\tau_i=\tau_i-C_\Sigma\), maximize
\[
g_i(r)=r\exp\{-a[r^2/2-\widetilde\tau_i]_+\},\qquad r\ge0.
\]
If \(\widetilde\tau_i<0\), the taper is active at zero and the maximum
occurs at \(r=a^{-1/2}\), giving
\(a^{-1/2}e^{a\widetilde\tau_i-1/2}\). If
\(\widetilde\tau_i\ge0\), the untapered plateau ends at
\(r_\tau=\sqrt{2\widetilde\tau_i}\). When
\(a^{-1/2}\le r_\tau\), the maximum is \(r_\tau\); otherwise it is the
same interior-tail value. These cases show that \(g_i\) has a finite
maximum, and the triangle inequality gives \(B_\omega<\infty\).

For near-field fidelity, let \(\mu^*\) be an uncontrolled equilibrium
with \(D_{\mu^*}(a_i)<\tau_i\) for every negative atom. Strictness and
continuity provide a neighborhood in which every \(\omega_i=1\).
The controlled and uncontrolled fields, including their Jacobians, are
therefore identical there.

Finally, along any sequence \(\|\mu_n\|_2\to\infty\), each fixed atom
satisfies \(D_{\mu_n}(a_i)\to\infty\), hence
\(\omega_i(D_{\mu_n}(a_i))\to0\). Finiteness of the atom set gives
\(\rho_{\mathrm{eff}}(\mu_n)\to0\).
\end{proof}

\subsection{Derivation and Proof of
Theorem~\ref{thm:aggregate_equilibria}}
\label{app:aggregate_equilibrium}

At a fixed context, use a regular minimal exponential-family
representation
\[
\pi_\xi(a)
=
h(a)\exp\{\xi^\top T(a)-\psi(\xi)\},
\qquad
\xi\in\mathcal H,
\]
and let
\[
\mathcal M
=
\{\nabla\psi(\xi):\xi\in\mathcal H\}
\]
be the attainable mean-parameter space. For a fixed-covariance Gaussian,
\(\xi=\Sigma^{-1}\mu\), \(T(a)=a\), and
\(\nabla\psi(\xi)=\mu\). For a categorical policy, a reference-logit
gauge gives the corresponding minimal \((K-1)\)-dimensional
representation. Regularity and minimality imply
\(\nabla^2\psi(\xi)\succ0\) on \(\mathcal H\) and make
\(\nabla\psi:\mathcal H\to\mathcal M\) one-to-one.

Write
\[
\widehat A=\widehat A^+-\widehat A^-,
\qquad
p=\mathbb E_\nu[\widehat A^+],
\qquad
q=\mathbb E_\nu[\widehat A^-],
\]
and define the unnormalized moments
\[
\bar{\mathbf m}_+
=
\mathbb E_\nu[\widehat A^+T(a)],
\qquad
\bar{\mathbf m}_-
=
\mathbb E_\nu[\widehat A^-T(a)].
\]
Because \(p>0\), let
\(\mathbf m_+=\bar{\mathbf m}_+/p\); when \(q>0\), let
\(\mathbf m_-=\bar{\mathbf m}_-/q\). If \(q=0\), use only
\(\bar{\mathbf m}_-=\mathbf0\).

Set
\[
b
=
\bar{\mathbf m}_+-\bar{\mathbf m}_-.
\]
Up to a constant independent of \(\xi\), the signed objective, field,
and Hessian are
\[
\mathcal L(\xi)
=
b^\top\xi-(p-q)\psi(\xi)+C,
\qquad
\mathbf F(\xi)
=
b-(p-q)\nabla\psi(\xi),
\]
\[
J_{\mathbf F}(\xi)
=
\nabla^2\mathcal L(\xi)
=
-(p-q)\nabla^2\psi(\xi).
\]
We consider
\[
\dot\xi=\mathbf F(\xi),
\qquad
\xi_{t+1}=\xi_t+\alpha\mathbf F(\xi_t),
\qquad
\alpha>0.
\]

\begin{proof}
\textbf{Positive-dominant regime, including positive-only
\((p>q\ge0)\).}
Define
\[
\mathbf m^\star
=
\frac{b}{p-q}.
\]
A finite equilibrium is equivalent to
\[
\nabla\psi(\xi^\star)=\mathbf m^\star.
\]
Thus, if \(\mathbf m^\star\in\mathcal M\), existence and uniqueness
follow from the bijectivity of
\(\nabla\psi:\mathcal H\to\mathcal M\); if
\(\mathbf m^\star\notin\mathcal M\), no finite equilibrium exists.
Because \(p-q>0\),
\[
\nabla^2\mathcal L(\xi)
=
-(p-q)\nabla^2\psi(\xi)
\prec0,
\]
so any finite equilibrium is the unique finite maximizer of
\(\mathcal L\). Its continuous-time Jacobian is negative definite, hence
it is locally asymptotically stable. The discrete local Jacobian is
\[
J_{\mathrm{disc}}(\xi^\star)
=
I-\alpha(p-q)\nabla^2\psi(\xi^\star),
\]
whose eigenvalues lie strictly inside the unit disk whenever
\[
0
<
\alpha
<
\frac{2}{
(p-q)\lambda_{\max}(\nabla^2\psi(\xi^\star))
}.
\]

When \(q=0\),
\(\bar{\mathbf m}_-=\mathbf0\) and
\(\mathbf m^\star=\mathbf m_+\), recovering the positive-only case.
When \(q>0\), with \(\rho=q/p\),
\[
\mathbf m^\star
=
\frac{p\mathbf m_+-q\mathbf m_-}{p-q}
=
\mathbf m_+
+
\frac{\rho}{1-\rho}
(\mathbf m_+-\mathbf m_-),
\]
so
\[
\mathbf m^\star-\mathbf m_+
=
\frac{\rho}{1-\rho}
(\mathbf m_+-\mathbf m_-).
\]
For fixed unequal \(\mathbf m_+\) and \(\mathbf m_-\), this algebraic
signed-target displacement grows without bound as \(\rho\to1^-\); a
finite equilibrium exists only while the target remains in
\(\mathcal M\).

If a sequence of reachable targets
\(\mathbf m_n^\star\in\mathcal M\) leaves every compact subset of
\(\mathcal M\), then the associated
\(\xi_n^\star=(\nabla\psi)^{-1}(\mathbf m_n^\star)\) leaves every compact
subset of \(\mathcal H\). Otherwise some compact
\(K\subset\mathcal H\) would contain a subsequence of
\(\xi_n^\star\), while its continuous image \(\nabla\psi(K)\) would be
compact in \(\mathcal M\) and contain the corresponding targets, a
contradiction.

\medskip
\noindent\textbf{Critical regime \((p=q>0)\).}
Here
\[
\mathbf F(\xi)
=
\bar{\mathbf m}_+-\bar{\mathbf m}_-
=
p(\mathbf m_+-\mathbf m_-)
\]
is constant. If \(\mathbf m_+\ne\mathbf m_-\),
\[
\xi(t)
=
\xi(0)+tp(\mathbf m_+-\mathbf m_-),
\qquad
\xi_t
=
\xi_0+t\alpha p(\mathbf m_+-\mathbf m_-),
\]
so both dynamics exhibit linear drift and have no finite equilibrium;
the drift persists for all forward time when \(\mathcal H\) is full or
forward invariant along this direction. If
\(\mathbf m_+=\mathbf m_-\), the field vanishes identically. Every point
is stationary, but no point is an isolated locally asymptotically stable
equilibrium.

\medskip
\noindent\textbf{Negative-dominant regime \((p<q)\).}
Now
\[
\nabla^2\mathcal L(\xi)
=
(q-p)\nabla^2\psi(\xi)
\succ0.
\]
Thus \(\mathcal L\) is strictly convex. If a finite stationary point
\(\xi^\star\) exists, its continuous field Jacobian is positive
definite, so it is repelling. The discrete local Jacobian is
\[
I+\alpha(q-p)\nabla^2\psi(\xi^\star),
\]
whose eigenvalues are
\(1+\alpha(q-p)\lambda_i>1\) for every \(\alpha>0\).
Hence every finite stationary point is unstable in both time forms; if
no stationary point exists, there is no finite stable equilibrium.
\end{proof}

\subsection{Proof of
Theorem~\ref{thm:policy_family_manifestations}}
\label{app:policy_family_manifestations}

\begin{proof}
Let \(r=q-p>0\).

\medskip
\noindent\textbf{Gaussian policy.}
Define the unnormalized signed moments
\[
\bar\mu_+
=
\mathbb E_\nu[\widehat A^+(a)a],
\qquad
\bar\mu_-
=
\mathbb E_\nu[\widehat A^-(a)a].
\]
For a fixed-covariance Gaussian,
\[
\nabla_\mu\log\pi_\mu(a)
=
\Sigma^{-1}(a-\mu),
\]
so
\[
\mathbf F_\mu(\mu)
=
\Sigma^{-1}
\bigl[\bar\mu_+-\bar\mu_--(p-q)\mu\bigr].
\]
Define
\[
\mu^\dagger
=
\frac{\bar\mu_--\bar\mu_+}{r}.
\]
Then
\[
\mathbf F_\mu(\mu)
=
r\Sigma^{-1}(\mu-\mu^\dagger).
\]
Writing \(\delta=\mu-\mu^\dagger\), the continuous and discrete
solutions are
\[
\delta(t)
=
e^{r\Sigma^{-1}t}\delta(0),
\qquad
\delta_t
=
\left(I+\alpha r\Sigma^{-1}\right)^t\delta_0.
\]
Since \(\Sigma^{-1}\succ0\),
\[
\|\delta(t)\|_2
\ge
e^{r\lambda_{\min}(\Sigma^{-1})t}\|\delta(0)\|_2,
\]
and
\[
\|\delta_t\|_2
\ge
\left[
1+\alpha r\lambda_{\min}(\Sigma^{-1})
\right]^t
\|\delta_0\|_2.
\]
Thus \(\mu^\dagger\) is repelling and every trajectory not initialized
exactly there diverges. For any fixed historical action \(a\),
\[
\|\nabla_\mu\log\pi_\mu(a)\|_2
=
\|\Sigma^{-1}(a-\mu)\|_2
\ge
\lambda_{\min}(\Sigma^{-1})\|a-\mu\|_2
\longrightarrow\infty.
\]

\medskip
\noindent\textbf{Categorical policy.}
Work on the zero-sum gauge
\[
\mathcal H_0
=
\{z\in\mathbb R^K:\mathbf1^\top z=0\}.
\]
Let the fixed signed action-mass vector be
\[
b
=
\mathbb E_\nu[\widehat A(a)e_a]
=
\mathbb E_\nu[\widehat A^+(a)e_a]
-
\mathbb E_\nu[\widehat A^-(a)e_a].
\]
Then
\[
\mathcal L(z)
=
b^\top z
+
r\log\sum_{j=1}^K e^{z_j}
+
C,
\]
with gradient and Hessian
\[
g(z)
=
\nabla\mathcal L(z)
=
b+r\pi_z,
\qquad
H(z)
=
r\bigl(\operatorname{Diag}(\pi_z)-\pi_z\pi_z^\top\bigr).
\]
Here \(g\) is the aggregate objective gradient, not the single-action
score whose norm is bounded in Proposition~\ref{prop:distance_strength}.
Because
\(\mathbf1^\top b=p-q=-r\),
\(\mathbf1^\top g(z)=0\), so both continuous and discrete dynamics remain
in \(\mathcal H_0\). For every nonzero \(v\in\mathcal H_0\),
\[
v^\top H(z)v
=
r\,\operatorname{Var}_{A\sim\pi_z}(v_A)
>
0.
\]
Hence \(\mathcal L\) is strictly convex on \(\mathcal H_0\), has at most
one finite stationary point, and any such point is repelling: the
continuous Jacobian restricted to \(\mathcal H_0\) is positive
definite, while the discrete Jacobian restricted to \(\mathcal H_0\)
has eigenvalues strictly larger than one for every \(\alpha>0\).

It remains to prove escape for every nonstationary initialization. If
\(g(z_0)=0\), both the continuous and discrete dynamics remain at the
stationary point. Otherwise, since \(g(z)=b+r\pi_z\) is bounded, the
continuous flow \(\dot z=g(z)\) exists for all forward time, and
\[
c_0=\|g(z_0)\|_2^2>0.
\]
Along the flow,
\[
\frac{d}{dt}\|g(z(t))\|_2^2
=
2g(z(t))^\top H(z(t))g(z(t))
\ge0,
\]
so \(\|g(z(t))\|_2^2\ge c_0\), and therefore
\[
\frac{d}{dt}\mathcal L(z(t))
=
\|g(z(t))\|_2^2
\ge c_0.
\]
Thus \(\mathcal L(z(t))\ge\mathcal L(z_0)+c_0t\).

For the discrete update
\[
z_{t+1}=z_t+\alpha g_t,
\qquad
g_t=g(z_t),
\]
gradient monotonicity gives
\[
(g_{t+1}-g_t)^\top(z_{t+1}-z_t)\ge0.
\]
Since \(z_{t+1}-z_t=\alpha g_t\),
\(g_{t+1}^\top g_t\ge\|g_t\|_2^2\), hence
\(\|g_{t+1}\|_2\ge\|g_t\|_2\ge\sqrt{c_0}\). Convexity then gives
\[
\mathcal L(z_{t+1})-\mathcal L(z_t)
\ge
g_t^\top(z_{t+1}-z_t)
=
\alpha\|g_t\|_2^2
\ge
\alpha c_0.
\]
Thus the discrete objective also grows at least linearly.

A continuous function is bounded above on every compact subset of
\(\mathcal H_0\). The preceding linear-growth bounds therefore imply
that every nonstationary continuous or discrete trajectory eventually
leaves every compact subset of \(\mathcal H_0\).

Finally, on \(\mathcal H_0\),
\[
z_i
=
\log\pi_z(i)
-
\frac1K\sum_{j=1}^K\log\pi_z(j).
\]
If the probability vectors stayed in a compact subset of the simplex
interior, all coordinates would be bounded below by some
\(\varepsilon>0\), and the inverse-softmax formula would keep the logits
bounded. Escape in \(\mathcal H_0\) therefore implies a subsequence along
which at least one action probability approaches zero. Nevertheless,
for every action \(a\),
\[
\|\nabla_z\log\pi_z(a)\|_2^2
=
(1-\pi_z(a))^2+\sum_{j\ne a}\pi_z(j)^2
\le
2(1-\pi_z(a))^2
\le2,
\]
as established in Proposition~\ref{prop:distance_strength}.
\end{proof}

\subsection{Proof of Minimum-Order Control}
\label{app:min_order_control}

For fixed-covariance Gaussian mean coordinates,
Proposition~\ref{prop:distance_strength} gives
\(R(D)=\Theta(D-C_\Sigma)\). Since \(C_\Sigma\) is fixed,
\(D-C_\Sigma=\Theta(D)\) as \(D\to\infty\). Hence there are
constants \(0<c_1\le c_2<\infty\) and \(D_0\) such that, for all
\(D\ge D_0\),
\[
c_1D
\le
R(D)
\le
c_2D.
\]
For the nonnegative weight \(\omega(D)\),
\[
c_1D\,\omega(D)
\le
\omega(D)R(D)
\le
c_2D\,\omega(D).
\]
Thus \(\omega(D)R(D)\) is bounded if and only if
\(\omega(D)=O(D^{-1})\), and it vanishes if and only if
\(\omega(D)=o(D^{-1})\). Since
\(D-C_\Sigma=r^2/2\), the matching order is equivalently
\(O(r^{-2})\) in standardized distance. This statement concerns the
squared-score reuse loop gain \(R\); the parameter-force scale is
\(\sqrt R\), as distinguished in the main text.

\subsection{Empirical Negative-Update Utility}
\label{app:negative_update_utility}

Let \(\widehat J_{\mathcal B}(\theta)\) be an empirical task-performance
estimate on evaluation batch \(\mathcal B\). For a negative sample
\(z=(s,a)\), define
\[
g_\theta^-(z)
=
-\widehat A^-(z)\nabla_\theta\log\pi_\theta(a\mid s),
\qquad
\Delta\theta(z;w)
=
\eta w g_\theta^-(z),
\qquad
0\le w\le1.
\]
Its empirical and population utilities are
\[
\widehat{\mathcal U}_{\mathcal B}(z;w\mid\theta)
=
\widehat J_{\mathcal B}\bigl(\theta+\Delta\theta(z;w)\bigr)
-
\widehat J_{\mathcal B}(\theta),
\]
\[
\mathcal U_\theta(z;w)
=
J\bigl(\theta+\eta w g_\theta^-(z)\bigr)-J(\theta).
\]
The empirical utility-optimal weight is
\[
\widehat w_{\mathcal U}^\star(z\mid\theta)
\in
\arg\max_{w\in[0,1]}
\widehat{\mathcal U}_{\mathcal B}(z;w\mid\theta).
\]
If \(\widehat J_{\mathcal B}\) is unbiased for \(J\), then for fixed
\(\theta,z,w\),
\[
\mathbb E_{\mathcal B}
\left[
\widehat{\mathcal U}_{\mathcal B}(z;w\mid\theta)
\right]
=
\mathcal U_\theta(z;w).
\]
Thus empirical utility is evaluation-batch dependent rather than an
intrinsic deterministic sample label.

\begin{lemma}[Far-field shrinkage of the utility-optimal step]
\label{lem:far_field_utility}
Let \(D\) index learner-relative remoteness and let \(g(D)\ne0\) be the
corresponding negative-update direction. Define
\[
I(D)=\|g(D)\|_2,
\qquad
v(D)=\frac{g(D)}{I(D)}.
\]
Suppose the local update space admits an
orthogonal decomposition
\[
\mathcal T
=
\mathcal T_{\mathrm{task}}
\oplus
\mathcal T_{\mathrm{far}},
\]
with projection \(P_{\mathrm{task}}\), and suppose that:
\begin{enumerate}
    \item \(\nabla J(\theta)\in\mathcal T_{\mathrm{task}}\);

    \item \(\|P_{\mathrm{task}}v(D)\|_2\to0\) as \(D\to\infty\);

    \item there exist \(\rho_0>0\), \(\kappa>0\), and \(D_0\) such that,
    for every \(D\ge D_0\) and every \(\rho\in[0,\rho_0]\),
    \[
    v(D)^\top
    \nabla^2J\bigl(\theta+\rho v(D)\bigr)
    v(D)
    \le
    -\kappa.
    \]
\end{enumerate}
Define
\[
\mathcal V_D(\rho)
=
J\bigl(\theta+\rho v(D)\bigr)-J(\theta),
\qquad
0\le\rho\le\rho_0.
\]
Continuity of \(J\) implies that \(\mathcal V_D\) attains a maximum
on the compact interval \([0,\rho_0]\); let
\[
\rho_{\mathcal U}^\star(D)
\in
\arg\max_{\rho\in[0,\rho_0]}
\mathcal V_D(\rho).
\]
Then
\[
\rho_{\mathcal U}^\star(D)\to0
\qquad
\text{as }D\to\infty.
\]
Moreover, for \(D\ge D_0\), if
\(\langle\nabla J(\theta),v(D)\rangle<0\), then
\[
\mathcal V_D(\rho)<0
\qquad
\text{for every }\rho\in(0,\rho_0].
\]
\end{lemma}

\begin{proof}
Let
\[
a(D)
=
\langle\nabla J(\theta),v(D)\rangle.
\]
Because
\(\nabla J(\theta)\in\mathcal T_{\mathrm{task}}\),
\[
|a(D)|
=
\left|
\left\langle
\nabla J(\theta),
P_{\mathrm{task}}v(D)
\right\rangle
\right|
\le
\|\nabla J(\theta)\|_2
\|P_{\mathrm{task}}v(D)\|_2
\longrightarrow0.
\]
For \(D\ge D_0\) and \(\rho\in(0,\rho_0]\), Taylor's theorem gives some
\(\xi\in(0,\rho)\), and the curvature assumption yields
\[
\mathcal V_D(\rho)
\le
\rho a(D)-\frac{\kappa}{2}\rho^2.
\]
The same bound is trivial at \(\rho=0\). Since
\(\mathcal V_D(0)=0\), a maximizing step must have nonnegative utility.
The right-hand side is strictly negative whenever
\[
\rho
>
\frac{2[a(D)]_+}{\kappa},
\]
so
\[
0
\le
\rho_{\mathcal U}^\star(D)
\le
\frac{2[a(D)]_+}{\kappa}
\longrightarrow0.
\]
If \(a(D)<0\), the same upper bound is strictly negative for every
\(\rho>0\), proving the second claim.
\end{proof}

This is a sufficient local mechanism, not a universal monotonicity law.
It applies when far-field negative updates become asymptotically
misaligned with the task-improving direction and the evaluated objective
is locally harmed along those directions. The curvature condition can
represent local losses caused by parameter drift or policy-support
contraction, provided those effects reduce the evaluated objective
\(J\). If far-field growth remains aligned with \(\nabla J(\theta)\),
the conclusion need not hold.

A sample weight \(w\) produces effective step length
\[
\rho(D;w)
=
\eta w I(D).
\]
Because the lemma is local to \([0,\rho_0]\), define
\[
w_{\mathcal U,\mathrm{loc}}^\star(D)
\in
\arg\max_{\substack{0\le w\le1\\ \eta wI(D)\le\rho_0}}
\mathcal V_D\bigl(\eta wI(D)\bigr).
\]
The same quadratic bound and feasibility of the zero step yield
\[
0
\le
\eta I(D)w_{\mathcal U,\mathrm{loc}}^\star(D)
\le
\frac{2[a(D)]_+}{\kappa}
\longrightarrow0.
\]
For fixed \(\eta>0\), if
\(\liminf_{D\to\infty}I(D)>0\), then
\(w_{\mathcal U,\mathrm{loc}}^\star(D)\to0\) as well.

\subsection{Proof of Proposition~\ref{prop:exp_vanishing_influence}}
\label{app:exp_vanishing_influence}

\begin{proof}
For \(D>\tau\),
\[
\omega_{\mathrm{Exp}}(D)
=
\exp\left[-\frac{\lambda}{c}(D-\tau)\right].
\]
Under the finite-order response assumption,
\[
0
\le
\omega_{\mathrm{Exp}}(D)I(D)
\le
C e^{(\lambda/c)\tau}
e^{-(\lambda/c)D}(1+D)^m
\longrightarrow0,
\]
because exponential decay dominates every finite-order polynomial.

When the response scale is instantiated by a raw negative-update
vector, let \(z_D\) denote a corresponding sample and write
\[
I(D)=\|g_\theta^-(z_D)\|_2.
\]
Define
\[
\Delta\theta_{\mathrm{Exp}}(D)
:=
\eta\omega_{\mathrm{Exp}}(D)g_\theta^-(z_D).
\]
Then
\[
\|\Delta\theta_{\mathrm{Exp}}(D)\|_2
=
\eta\omega_{\mathrm{Exp}}(D)I(D)
\longrightarrow0.
\]
If \(\widehat J_{\mathcal B}\) is locally
\(L_{\mathcal B}\)-Lipschitz around \(\theta\), then for all sufficiently
large \(D\),
\[
\left|
\widehat{\mathcal U}_{\mathcal B}
\bigl(z_D;\omega_{\mathrm{Exp}}(D)\mid\theta\bigr)
\right|
\le
L_{\mathcal B}
\|\Delta\theta_{\mathrm{Exp}}(D)\|_2
\longrightarrow0.
\]
This proves the proposition together with the parameter-step and
conditional empirical-utility consequences.
\end{proof}

\paragraph{Detailed guarantee hierarchy.}
The control conditions form a guarantee hierarchy. L0,
\(\omega(D)\to0\), gives eventual inward Gaussian radial drift for each
fixed finite \(\rho\); L1, \(\omega(D)\sqrt D\to0\), removes the
absolute negative far-field force; L2, \(\omega(D)D=O(1)\), bounds the
Gaussian reuse loop gain; and L3,
\(D^k\omega(D)\to0\) for every finite \(k\), removes all finite-order
tails. Global scaling fails L0, reciprocal-quadratic tapering is critical at
L2, and exponential tapering satisfies L0--L3.

\paragraph{Control mapping.}
Positive-only and global scaling ignore remoteness; reciprocal tapers have
polynomial tails, and the quadratic member meets the matching order; DRPO uses
the exponential member that satisfies L0--L3. These controls share the actor
form in Equation~\eqref{eq:drpo_actor_field} and differ only in their
negative-branch coefficient
\citep{peng2019advantage,nair2020awac,kostrikov2021offline,
grigsby2021closer,arnal2025asymmetric,leroux2025tapered}.


\subsection{Categorical One-vs-Rest Dynamics}
\label{app:categorical_self_limiting}


\begin{theorem}[Categorical self-limiting suppression]
\label{thm:categorical_self_limiting}
For a repeatedly suppressed action \(i\), define
\(y_i=z_i-\log\sum_{j\ne i}e^{z_j}\), so \(\pi_i=\sigma(y_i)\).
Under an uncontrolled negative atom of mass \(q>0\),
\[
\dot y_i=-q(1-\pi_i)c_K(t),
\qquad \frac K{K-1}\le c_K(t)\le2.
\]
Hence \(y_i\) eventually decreases linearly and \(\pi_i\) exponentially.
With \(\omega_i=\min\{1,(e^{\tau_i}\pi_i)^\beta\}\),
\(\beta=\lambda/c\), the far branch instead satisfies
\[
y_i(t)=-\beta^{-1}\log t+O(1),
\qquad \pi_i(t)=\Theta(t^{-1/\beta}).
\]
Thus tapering changes persistent suppression from exponential to
polynomial probability decay; it does not freeze the action probability.
\end{theorem}

The finite-crossing and restoring-force conditions are recorded separately in
Appendix~\ref{app:categorical_restoring_balance}.

\begin{proof}[Proof of Theorem~\ref{thm:categorical_self_limiting}]
For a single negative atom \(i\), the direct-logit field is
\[
\dot z=-q(e_i-\pi),
\qquad
\dot z_i=-q(1-\pi_i),
\qquad
\dot z_j=q\pi_j\quad(j\ne i).
\]
Define the normalized off-target probabilities
\(\bar\pi_j=\pi_j/(1-\pi_i)\). Differentiating
\(y_i=z_i-\log\sum_{j\ne i}e^{z_j}\) gives
\begin{align}
\dot y_i
&=-q(1-\pi_i)
-q\frac{\sum_{j\ne i}\pi_j^2}{1-\pi_i} \\
&=-q(1-\pi_i)c_K(t),
\qquad
c_K(t)=1+\sum_{j\ne i}\bar\pi_j^2.
\label{eq:app_categorical_exact_ode}
\end{align}
Because \(\bar\pi\) is a distribution on \(K-1\) actions,
\[
\frac1{K-1}\le\sum_{j\ne i}\bar\pi_j^2\le1,
\]
which proves the coefficient envelope without a concentration assumption.

We next justify entry into the far half-line before using its bounds.
The log-odds is strictly decreasing. Until \(\pi_i<1/2\), monotonicity
gives \(1-\pi_i(t)\ge1-\pi_i(0)>0\), so
Equation~\eqref{eq:app_categorical_exact_ode} has a strictly negative
upper bound. The trajectory therefore crosses \(\pi_i=1/2\) in finite
time. Thereafter \(1-\pi_i\ge1/2\), so \(y_i\to-\infty\) at a linear
rate and \(\pi_i=\sigma(y_i)\) decays exponentially.

For the generic coefficient, suppose \(j^*\ne i\) is the unique largest
off-target logit initially. For every \(k\ne i,j^*\),
\[
\frac{d}{dt}(z_{j^*}-z_k)=q(\pi_{j^*}-\pi_k)>0.
\]
The initial ordering is therefore preserved. If a gap stayed bounded,
then, once \(\pi_i\to0\), maximality of \(j^*\) and the positive initial
gap would give a positive lower bound on
\(\pi_{j^*}-\pi_k\), a contradiction. Hence all such gaps diverge,
off-target mass concentrates on \(j^*\), and \(c_K(t)\to2\). Under
complete off-target symmetry, symmetry is preserved and
\(\bar\pi_j=1/(K-1)\), giving
\(c_K(t)\equiv K/(K-1)\).

Now apply the taper
\(\omega_i=\min\{1,(e^{\tau_i}\pi_i)^\beta\}\). Since the tapered
log-odds is strictly decreasing with a rate bounded away from zero on every
bounded log-odds interval, it reaches the far branch
\(\pi_i<\min\{1/2,e^{-\tau_i}\}\) in finite time. Thereafter let
\(U=e^{-\beta y_i}\). Since
\(\pi_i e^{-y_i}=1-\pi_i\),
\begin{align}
\dot U
&=-\beta e^{-\beta y_i}\dot y_i \\
&=\beta q e^{\beta\tau_i}
c_K(t)(1-\pi_i)^{\beta+1}.
\label{eq:app_categorical_tapered_u}
\end{align}
Thereafter, the right-hand side is bounded above
and below by positive constants. Thus \(U=\Theta(t)\), which is
equivalent to
\[
y_i(t)=-\beta^{-1}\log t+O(1),
\qquad
\pi_i(t)=\Theta(t^{-1/\beta}).
\]
\end{proof}

\subsection{Categorical Restoring-Force Balance}
\label{app:categorical_restoring_balance}

Write the one-dimensional restored dynamics as
\[
\dot y=F(y)=E(y)-S_{\mathrm{tap}}(y),
\]
where \(S_{\mathrm{tap}}>0\) is downward suppression and \(E>0\) is
upward restoration.

\begin{proposition}[Finite restoring crossing]
\label{prop:categorical_restoring_crossing}
Assume \(E\) and \(S_{\mathrm{tap}}\) are continuous and the zeros of
\(F=E-S_{\mathrm{tap}}\) in \([y_L,y_R]\) are isolated. If
\(F(y_L)>0\) and \(F(y_R)<0\), then the interval contains a zero at
which the phase line has a positive-to-negative crossing. Such a
crossing is locally asymptotically stable; if it is hyperbolic, its
condition is
\[
F'(y^*)<0
\quad\Longleftrightarrow\quad
[S_{\mathrm{tap}}-E]'(y^*)>0.
\]
\end{proposition}

\begin{proof}
The intermediate value theorem gives at least one zero. Because the
zeros of \(F\) are isolated on a compact interval, only finitely many occur; as
the endpoint signs differ, at least one has positive sign immediately
to its left and negative sign immediately to its right. The flow then
points toward that zero from both sides, proving phase-line stability. At a
hyperbolic crossing this sign change is equivalent to \(F'(y^*)<0\),
and the displayed derivative identity follows from \(F=E-S_{\mathrm{tap}}\).
\end{proof}

In the binary or symmetric one-vs-rest tail, the tapered suppression
has order \(e^{\beta y}\), whereas entropy restoration has order
\(|y|e^y\) as \(y\to-\infty\). Hence \(\beta\ge1\) supplies the
required far-tail ordering (with the \(|y|\) factor breaking the tie at
\(\beta=1\)), but it does not by itself verify the finite phase-line
conditions of Proposition~\ref{prop:categorical_restoring_crossing}.

\section{Distributional Reweighting Interpretations}
\label{app:distributional_reweighting}

\subsection{Finite-Measure I-Divergence Attenuation}
\label{app:finite_measure_variational}

DRPO's exponential attenuation admits an optimistic finite-measure
reweighting interpretation analogous to divergence-ball DRO, with the
constraint expressed in generalized I-divergence
\citep{csiszar1975divergence}. Optimistic or best-case distributional
optimization also has a formal precedent in likelihood approximation
and policy optimization
\citep{nguyen2019optimistic,song2020optimistic}. Our construction is
related but distinct: it reweights a finite measure of negative-update
mass according to learner-relative remoteness and may reduce that
measure's total mass. It is therefore an analogue of optimistic
divergence-ball reweighting, not a standard ambiguity set over normalized
probability distributions. The analogy is deliberately qualified: unlike
standard ambiguity sets over normalized probability measures, the present
construction permits attenuation of the total negative-update mass. Let
\(d\xi^-(z)=\widehat A^-(z)d\nu(z)\) and
\(h(z)=[D(z)-\tau(z)]_+\). Consider
\begin{equation}
\min_{0\le w\le1}\int wh\,d\xi^-
\quad\mathrm{s.t.}\quad
\int(w\log w-w+1)d\xi^-\le\delta.
\label{eq:finite_measure_variational}
\end{equation}

\begin{proposition}[Finite-measure exponential reweighting]
\label{prop:finite_measure_variational}
Let \(\delta_0=\xi^-(\{h>0\})\). For
\(0<\delta<\delta_0\), the constraint in
Equation~\eqref{eq:finite_measure_variational} is active and has a unique
finite \(\kappa>0\) with
\(w_\kappa(z)=e^{-h(z)/\kappa}\). At \(\delta=\delta_0\), the endpoint is
\(w_0=\mathbf1\{h=0\}\); for \(\delta>\delta_0\), zero objective is
already attainable and no finite one-to-one \(\kappa\)--\(\delta\)
relation remains. The later endpoint \(\delta\ge|\xi^-|\) merely makes
\(w\equiv0\) feasible.
\end{proposition}

\begin{proof}[Proof of Proposition~\ref{prop:finite_measure_variational}]
Use the convention
\(\phi(0)=1\) for
\(\phi(w)=w\log w-w+1\). For an active constraint, the Lagrangian with
multiplier \(\kappa>0\) is
\[
\mathcal L(w,\kappa)
=\int\{wh+\kappa\phi(w)\}\,d\xi^- -\kappa\delta.
\]
Pointwise stationarity on \(0<w\le1\) gives
\(h+\kappa\log w=0\), hence
\[
w_\kappa(z)=e^{-h(z)/\kappa}.
\]
For \(h=0\), this equals one; for \(h>0\), it lies strictly between
zero and one. Define
\(\mathcal D(\kappa)=\int\phi(w_\kappa)d\xi^-\). Dominated convergence gives
\[
\lim_{\kappa\to\infty}\mathcal D(\kappa)=0,
\qquad
\lim_{\kappa\downarrow0}\mathcal D(\kappa)
=\xi^-(\{h>0\})=\delta_0.
\]
On a set of positive \(h\), increasing \(\kappa\) strictly increases
\(w_\kappa\) toward one and strictly decreases \(\phi(w_\kappa)\).
Thus \(\mathcal D(\kappa)\) is continuous and strictly decreasing from
\(\delta_0\) to zero, proving existence and uniqueness for
\(0<\delta<\delta_0\).

At \(\delta=\delta_0\), the limit is
\(w_0=\mathbf1\{h=0\}\), which has zero objective and divergence
\(\delta_0\). For \(\delta>\delta_0\), the same zero-objective point is
strictly feasible, so the constraint need not be active and there is no
finite one-to-one multiplier relation. Finally,
\(\int\phi(0)d\xi^-=|\xi^-|\), so \(w\equiv0\) becomes feasible only at the later
endpoint \(\delta\ge|\xi^-|\).
\end{proof}

This problem derives learner-relative exponential attenuation. It is
distinct from quality-based hard selection: both are optimistic
reweighting views, but one controls remoteness under a finite negative
measure and the other selects a quality region. Neither problem implies
the other.

\subsection{Relation to Ratio- and Constraint-Based Controls}
\label{app:topr_relation}

\paragraph{Behavior-ratio identity.}
For total log probabilities, define
\[
r_{\mathrm{beh},\theta}(s,a)
=
\frac{\pi_\theta(a\mid s)}{\mu(a\mid s)},
\qquad
D_\theta(s,a)
=
-\log\pi_\theta(a\mid s),
\qquad
D_\mu(s,a)
=
-\log\mu(a\mid s),
\]
where \(\mu\) is the behavior policy. For any \(\beta>0\),
\[
\begin{aligned}
\min\left\{
    r_{\mathrm{beh},\theta}(s,a)^\beta,
    1
\right\}
&=
\min\left\{
    \left(
        \frac{\pi_\theta(a\mid s)}
             {\mu(a\mid s)}
    \right)^\beta,
    1
\right\} \\
&=
\exp\left\{
    -\beta
    \left[
        D_\theta(s,a)-D_\mu(s,a)
    \right]_+
\right\}.
\end{aligned}
\]
Thus, the clipped behavior ratio is exactly
Equation~\eqref{eq:drpo_exp_weight} with the sample-specific threshold
\(\tau=D_\mu(s,a)\) and rate \(\lambda/c=\beta\). At the atomic-action
level, canonical TOPR is the \(\beta=1\) case
\citep{leroux2025tapered}. The corresponding untruncated power weight
\(r_{\mathrm{beh},\theta}^{\beta}\) has the same far-field decay but may
exceed one when \(D_\theta<D_\mu\), whereas the thresholded taper remains
capped at one. Replacing both log probabilities by their
length-normalized values replaces
\(r_{\mathrm{beh},\theta}\) by
\(r_{\mathrm{beh},\theta}^{1/L}\).

For the fitted-reference comparator used in the structured-generation evaluation, the key distinction is the control coordinate: it uses the current-to-reference completion log-ratio, whereas DRPO uses learner-relative mean-token surprisal with a designed threshold and decay scale.

\paragraph{Conditional-versus-population constants.}

For a conditionally reused far-branch sample, the coefficient contains
\(\pi/\mu\), so the conditional vector field retains a factor of
\(\mu^{-1}\). In the population field, however, the outer sampling mass
gives
\[
\mu\min\{\pi/\mu,1\}
=
\min\{\pi,\mu\},
\]
and cancels the ratio denominator on the far branch.

\paragraph{Experimental fitted-reference variant.}
The TOPR arm implements the joint fitted-reference
\(\beta\)-TOPR variant; it is not a canonical frozen-behavior
reproduction of \citet{leroux2025tapered}. The canonical population
cancellation above requires the same \(\mu\) to serve as both the outer
sampling mass and the ratio denominator. It therefore does not apply to
the experimental variant, whose outer measure is induced by
prompt-uniform replay with a per-prompt unique-negative mean and whose
denominator is \(\mu_{\mathrm{ref}}\). The experimental arm also uses
full-completion summed log probabilities rather than the
length-normalized ratio identity above.

For Gaussian atoms, the triangle inequality and the one-atom \(G_i\)
calculation in Appendix~\ref{app:gaussian_drpo_boundedness} first give
weighted constants. A \(\mu_{\min}\) expression follows only after
imposing finite support and replacing the weighted terms by their
worst-case bounds.

\paragraph{Refreshed proximal ratios.}
For a PPO-style proximal update, define
\[
r_{\mathrm{prox},\theta}(s,a)
=
\frac{\pi_\theta(a\mid s)}
     {\pi_{\mathrm{old}}(a\mid s)}.
\]
For the clipped surrogate
\citep{schulman2017proximal}, when \(\widehat A<0\), the effective
coefficient relative to
\(\widehat A\nabla_\theta\log\pi_\theta(a\mid s)\), away from the clipping
boundary, is
\[
r_{\mathrm{prox},\theta}
\mathbf 1
\left\{
    r_{\mathrm{prox},\theta}
    \geq
    1-\epsilon
\right\},
\]
rather than a pure indicator. The activation boundary is hard, but the
retained coefficient remains the proximal policy ratio. Refreshing
\(\pi_{\mathrm{old}}\) resets this ratio to one at the beginning of each
outer iteration. The resulting gate therefore constrains movement relative
to the latest anchor but does not directly encode the cumulative
learner-relative remoteness that increases under
Theorem~\ref{thm:reuse_dynamics}.

\paragraph{Movement, tethering, and source control.}
Reducing a fixed positive step size does not restore a stable finite
equilibrium under negative dominance
(Theorem~\ref{thm:aggregate_equilibria}). Reference regularization instead
restores an inward force by tethering the whole policy. When \(q>p\),
stability requires
\[
\gamma>q-p.
\]
When \(p>q\), strengthening the tether contracts the equilibrium toward
the reference
(Proposition~\ref{prop:regularized_restoration}). DRPO instead attenuates
the learner-relative remote negative source while leaving the positive
branch and near-field negative feedback unchanged. These mechanisms are
not mutually exclusive, but the dynamics of their combination lie outside
the fixed-data, fixed-advantage analysis considered here.

\section{Experimental Details}
\label{app:experimental_details}

\subsection{D4RL Locomotion Gradient-Diagnostic Protocol}
\label{app:d4rl_gradient_diagnostic}

\paragraph{Critic fitting and frozen advantages.}
For each of the nine D4RL locomotion dataset cells, we train one canonical
value critic for 100,000 optimizer updates. We then compute, standardize, and
freeze the transition-level learned-critic advantages. All actor seeds for a
given dataset use the same frozen critic outputs.

\paragraph{Positive-only reference actors.}
For seeds 100--109, we train a Gaussian reference actor for 100,000 optimizer
updates using only transitions with positive frozen advantage. The terminal
actor checkpoint is fixed before the gradient diagnostic and serves only to
define learner-relative distance and the implemented-gradient measurement.

\paragraph{Matched negative-transition probe.}
The probe population contains transitions with negative frozen advantage.
Near and far pools use standardized-distance quantiles 0.25 and 0.75. We match
absolute advantage using 20 bins and a 5\% relative tolerance, retain 256
matched pairs per seed, and evaluate the full actor-parameter gradient on 64
pairs per seed. Seven far-distance bins are reported. The probe performs no
optimizer update.

\paragraph{Aggregation and provenance.}
Every far point is normalized by its original matched near partner. Curves are
first aggregated within seed and dataset; the main panel then assigns equal
weight to the nine dataset cells on their shared distance support. Its 95\%
confidence interval uses 10,000 hierarchical bootstrap draws, resampling
dataset cells and then seeds, with bootstrap seed 20260723. The archived
point-level output records dataset, critic, actor, RunSpec, source, and
checkpoint provenance.

\begin{table*}[t]
\centering
\caption{Configuration of the D4RL locomotion gradient diagnostic.}
\label{tab:app_d4rl_gradient_protocol}
\small
\setlength{\tabcolsep}{5pt}
\begin{tabular}{ll}
\toprule
Component & Configuration \\
\midrule
Datasets & HalfCheetah, Hopper, Walker2d $\times$ medium, medium-replay, medium-expert \\
Critic & One learned value critic per dataset; 100,000 optimizer updates \\
Advantage & Transition-level learned-critic advantage; standardized and frozen \\
Reference actor & Positive-only Gaussian actor; 100,000 optimizer updates \\
Seeds & 100--109 \\
Probe population & Negative frozen-advantage transitions \\
Near/far pools & Standardized-distance quantiles 0.25 / 0.75 \\
Advantage matching & 20 bins; 5\% relative tolerance \\
Matched pool & 256 pairs per seed \\
Full-gradient subset & 64 pairs per seed \\
Distance summary & Seven far-distance bins with original-pair normalization \\
Probe updates & None \\
Cross-dataset summary & Within-seed/dataset normalization; equal dataset-cell weighting \\
Uncertainty & Hierarchical bootstrap over dataset cells and seeds \\
\bottomrule
\end{tabular}
\end{table*}

\begin{figure*}[t]
\centering

\begin{minipage}[t]{0.49\textwidth}
    \centering
    \textbf{(a) D4RL}\\[-0.2em]
    \includegraphics[width=\linewidth]
    {figures/fig_app_d4rl9_gradient_panels.pdf}
\end{minipage}
\hfill
\begin{minipage}[t]{0.49\textwidth}
    \centering
    \textbf{(b) Structured-generation tasks}\\[-0.2em]
    \includegraphics[width=\linewidth]
    {figures/fig_app_structured9_gradient_panels.pdf}
\end{minipage}

\caption{
\textbf{Disaggregated external far-field diagnostics.}
(a) Nine D4RL locomotion dataset cells.
(b) Nine structured-generation tasks.
}
\Description{Side-by-side external diagnostic panels for the nine D4RL
dataset cells and nine structured-generation tasks.}
\label{fig:app_external_gradient_panels}
\end{figure*}

\subsection{Unified Continuous Environment C-U1}
\label{app:cu1_environment}

C-U1 is the shared controlled continuous environment used by all
continuous mechanism experiments. A context
$s\in\mathbb R^6$
is sampled from
$\mathcal N(0,I_6)$.
For each seed, we independently sample 4096 training contexts and 4096
test contexts from the same distribution. Test contexts never
participate in optimization and are used only to measure
held-out-context generalization.

\paragraph{State-dependent task geometry.}
Each context determines a two-dimensional positive-support center
$a_+(s)$ and a unit task direction $u(s)$. We define
\begin{equation}
\begin{aligned}
    [a_+(s)]_1
    &=
    0.70
    \tanh
    \bigl(
        0.85s_1
        -
        0.30s_2s_3
        +
        0.20\sin(1.6s_4)
    \bigr), \\
    [a_+(s)]_2
    &=
    0.65
    \tanh
    \bigl(
        -0.50s_2
        +
        0.35\cos(1.1s_5)
        +
        0.22s_1s_6
    \bigr).
\end{aligned}
\label{eq:cu1_positive_center}
\end{equation}
The direction angle is
\begin{equation}
\begin{aligned}
    \varphi(s)
    &=
    1.15
    \tanh
    \bigl(
        0.75s_1
        +
        0.50s_3
        -
        0.30s_6
    \bigr) \\
    &\quad+
    0.30\sin(1.35s_2),
\end{aligned}
\label{eq:cu1_direction_angle}
\end{equation}
with
\begin{equation}
    u(s)
    =
    \begin{bmatrix}
        \cos\varphi(s)\\
        \sin\varphi(s)
    \end{bmatrix},
    \qquad
    v(s)
    =
    \begin{bmatrix}
        -u_2(s)\\
        u_1(s)
    \end{bmatrix}.
    \label{eq:cu1_task_basis}
\end{equation}
The unobserved task optimum and the directionally useful negative anchor
are
\begin{equation}
\begin{aligned}
    a_\star(s)
    &=
    a_+(s)+0.70u(s), \\
    a_-(s)
    &=
    a_+(s)-0.50u(s).
\end{aligned}
\label{eq:cu1_task_points}
\end{equation}

\paragraph{Positive and negative actions.}
Each context is associated with four positive actions and eight negative
actions. The positive actions lie on a reward contour of radius
$r_+=0.75$
centered at $a_\star(s)$:
\begin{equation}
\begin{aligned}
    a_{\mathrm{pos},k}(s)
    =
    a_\star(s)
    +
    r_+
    \bigl(
        \cos\phi_k\,u(s)
        +
        \sin\phi_k\,v(s)
    \bigr),
\end{aligned}
\label{eq:cu1_positive_actions}
\end{equation}
where
\begin{equation}
\begin{aligned}
    \{\phi_k\}_{k=1}^{4}
    &=
    \{
        \pi-\theta_1,\,
        \pi+\theta_1,\,
        \pi-\theta_2,\,
        \pi+\theta_2
    \}, \\
    \theta_1
    &=
    0.20,
    \qquad
    \cos\theta_2
    =
    \frac{2(0.70)}{0.75}
    -
    \cos\theta_1.
\end{aligned}
\label{eq:cu1_positive_angles}
\end{equation}
This construction makes the four actions equal-reward while ensuring
that their centroid is exactly $a_+(s)$. It also leaves nonzero
conditional action spread, preventing deterministic maximum-likelihood
variance contraction from being confused with the far-field mechanism.

The eight negative actions lie on a second reward contour of radius
$r_-=1.20$:
\begin{equation}
\begin{aligned}
    a_{\mathrm{neg},j}(s)
    =
    a_\star(s)
    +
    r_-
    \bigl(
        \cos\psi_j\,u(s)
        +
        \sin\psi_j\,v(s)
    \bigr),
\end{aligned}
\label{eq:cu1_negative_actions}
\end{equation}
with
\begin{equation}
\begin{aligned}
    \{\psi_j\}_{j=1}^{8}
    =
    \left\{
        \pi,\,
        \frac{3\pi}{4},\,
        \frac{\pi}{2},\,
        \frac{\pi}{4},\,
        0,\,
        -\frac{\pi}{4},\,
        -\frac{\pi}{2},\,
        -\frac{3\pi}{4}
    \right\}.
\end{aligned}
\label{eq:cu1_negative_angles}
\end{equation}
The first negative action equals $a_-(s)$. All eight negative actions
have the same distance to $a_\star(s)$ and therefore the same reward and
fixed advantage, while having different distances from the current
policy.

\paragraph{Reward and fixed advantages.}
The ground-truth reward is
\begin{equation}
    R(s,a)
    =
    \exp
    \left(
        -
        \frac{
            \lVert a-a_\star(s)\rVert_2^2
        }{
            2(0.75)^2
        }
    \right).
    \label{eq:cu1_reward}
\end{equation}
Advantages are computed once using the fixed baseline $0.40$:
\begin{equation}
    \widehat A(s,a)
    =
    R(s,a)-0.40,
    \label{eq:cu1_advantage}
\end{equation}
and remain frozen throughout optimization. Consequently, all four
positive actions have positive advantages, all eight negative actions
have negative advantages, and negative-action quality is exactly matched
within each context.

\paragraph{Policy parameterization.}
The policy is an isotropic state-conditioned Gaussian,
\begin{equation}
    \pi_\theta(a\mid s)
    =
    \mathcal N
    \left(
        \mu_\theta(s),
        \sigma_\theta^2(s)I_2
    \right).
    \label{eq:cu1_policy}
\end{equation}
Both $\mu_\theta(s)$ and the scalar
$\log\sigma_\theta(s)$
are produced by a shared two-layer ReLU MLP with 64 hidden units per
layer and separate output heads. The initial standard deviation is
$0.60$, and the learnable-variance experiments use no artificial
variance clamp.  The continuous mechanism results reported in the main
text use the fixed-variance branch unless explicitly labeled otherwise. A finite crossing of
$\log\sigma_\theta(s)<-12$
is recorded as a support or variance-boundary event, whereas nonfinite
parameters or outputs are reported separately as NaN/Inf numerical
failure.

Training minibatches are sampled by context, and all actions associated
with a selected context are retrieved together. Thus, the replicated
actions within one context are not treated as independent contexts.
Positive and negative objectives are averaged separately before their
relative coefficient is applied, preventing the four-to-eight action
count ratio from mechanically changing the total negative-update mass.

\paragraph{Shared use across continuous protocols.}
All C-U1 experiments reuse the same contexts, reward geometry, fixed
advantages, and policy architecture. E1 compares equal-advantage
negative actions at different learner-relative distances. E2 applies
only positive updates and measures the same negative actions as phantom
probes. E3 dynamically partitions negative actions using the
standardized distance
\begin{equation}
    d_\theta(s,a)
    =
    \frac{
        \lVert a-\mu_\theta(s)\rVert_2
    }{
        \sigma_\theta(s)
    },
    \label{eq:cu1_standardized_distance}
\end{equation}
with the prespecified near--far threshold $d_\theta=5.0$, and applies
targeted near- and far-field interventions. E4 uses the aligned action
$a_-(s)$ as directionally useful local negative feedback and the
remaining contour actions as additional far-field pressure. The
experiment-specific coefficients, training horizons, seeds, taper
calibrations, and terminal criteria are reported separately with their
corresponding protocols.

\subsection{Controlled Categorical Environment D-U1}
\label{app:du1_environment}

D-U1 denotes the controlled categorical environment family used to
study persistent negative suppression, probability-support boundaries,
and shared-semantic generalization. Its E6 instantiation uses a fixed
catalogue of 64 unordered actions, each associated with a
four-dimensional unit semantic vector. Action identifiers are randomly
permuted for every seed, so numerical action indices carry no geometric
ordering.

For each seed, we independently sample 2048 training contexts and 2048
test contexts from
$s\sim\mathcal N(0,I_6)$.
The two splits follow the same context distribution, and test contexts
never participate in optimization. Results on this split are therefore
reported as held-out-context or unseen-state generalization rather than
OOD generalization.

\paragraph{Semantic action catalogue.}
Let
$\{e_i\}_{i=1}^{K}\subset\mathbb R^4$,
with $K=64$, denote the unit reward embeddings of the categorical
actions. They are sampled once per seed and randomly assigned to action
identifiers. For each context, two fixed linear maps generate a
demonstrated semantic target and a candidate improvement direction:
\begin{equation}
\begin{aligned}
    t_+(s)
    &=
    \operatorname{unit}
    \!\left(
        W_+^\top s
    \right), \\
    \widetilde u(s)
    &=
    W_u^\top s
    -
    \left\langle
        W_u^\top s,
        t_+(s)
    \right\rangle
    t_+(s), \\
    u(s)
    &=
    \operatorname{unit}
    \!\left(
        \widetilde u(s)
    \right).
\end{aligned}
\label{eq:du1_semantic_basis}
\end{equation}
The hidden optimal target and the local negative target are
\begin{equation}
\begin{aligned}
    t_\star(s)
    &=
    \operatorname{unit}
    \!\left(
        t_+(s)+\delta u(s)
    \right), \\
    t_-(s)
    &=
    \operatorname{unit}
    \!\left(
        t_+(s)-\delta u(s)
    \right),
\end{aligned}
\qquad
\delta=0.45.
\label{eq:du1_semantic_targets}
\end{equation}

\paragraph{Reward and action roles.}
The ground-truth reward of action $i$ is determined by its semantic
alignment with the hidden target:
\begin{equation}
    R(s,i)
    =
    \frac{1}{2}
    \left(
        1+
        t_\star(s)^\top e_i
    \right).
    \label{eq:du1_reward}
\end{equation}
The hidden optimal action is
\begin{equation}
    i_\star(s)
    =
    \arg\max_i
    t_\star(s)^\top e_i.
    \label{eq:du1_hidden_optimum}
\end{equation}
It is excluded from all positive demonstrations.

For each context, the four positive actions are the available actions
with the largest similarity to $t_+(s)$. After excluding the hidden
optimum and positive actions, the local negative action is the action
most aligned with $t_-(s)$. The four far-negative actions are selected
by their alignment with $-t_+(s)$. Thus,
\begin{equation}
\begin{aligned}
    \mathcal P(s)
    &=
    \operatorname{Top4}_{i\neq i_\star(s)}
    \,
    t_+(s)^\top e_i, \\
    i_{\mathrm{local}}(s)
    &=
    \arg\max_{i\notin
        \{i_\star(s)\}\cup\mathcal P(s)}
    t_-(s)^\top e_i, \\
    \mathcal F(s)
    &=
    \operatorname{Top4}_{i\notin
        \{i_\star(s),i_{\mathrm{local}}(s)\}
        \cup\mathcal P(s)}
    \,
    \bigl[-t_+(s)^\top e_i\bigr].
\end{aligned}
\label{eq:du1_action_roles}
\end{equation}
The positive and negative advantages are fixed throughout training:
\begin{equation}
    \widehat A(s,i)
    =
    \begin{cases}
        +1,
        & i\in\mathcal P(s),\\
        -1,
        & i=i_{\mathrm{local}}(s)
          \text{ or }i\in\mathcal F(s).
    \end{cases}
    \label{eq:du1_fixed_advantages}
\end{equation}
Consequently, differences between local and far-negative updates cannot
be attributed to different advantage magnitudes.

\paragraph{Shared-semantic categorical policy.}
The policy uses a two-layer tanh MLP with 64 hidden units per layer to
map the context to a unit semantic direction
$q_\theta(s)\in\mathbb R^4$.
Let $\widetilde e_i$ denote the action embedding used by the policy. The
categorical logits and probabilities are
\begin{equation}
\begin{aligned}
    \ell_{\theta,i}(s)
    &=
    \kappa_\theta(s)
    q_\theta(s)^\top
    \widetilde e_i, \\
    \pi_\theta(i\mid s)
    &=
    \frac{
        \exp\!\left(\ell_{\theta,i}(s)\right)
    }{
        \sum_{j=1}^{K}
        \exp\!\left(\ell_{\theta,j}(s)\right)
    }.
\end{aligned}
\label{eq:du1_policy}
\end{equation}
The concentration is either fixed at
$\kappa_\theta(s)=8$
or learned as
\begin{equation}
    \kappa_\theta(s)
    =
    \operatorname{softplus}
    \!\left(
        c_\theta(s)
    \right)
    +
    0.05,
    \label{eq:du1_concentration}
\end{equation}
with initial value 8 and no upper clamp.

In the semantically aligned condition,
$\widetilde e_i=e_i$.
In the shuffled control, the same policy embeddings are randomly
permuted across action identifiers while the reward embeddings and
training samples remain unchanged. This intervention preserves the
catalogue size and categorical parameterization but breaks the
correspondence between policy-side similarity and ground-truth semantic
utility.

\paragraph{Evaluation and failure events.}
We evaluate expected semantic reward, probability assigned to the hidden
optimal action, entropy, effective support, and normalized semantic
extrapolation on both training and held-out contexts. Task-performance
collapse, effective-support or concentration-boundary events, and
NaN/Inf numerical failure are reported separately. A policy can
therefore retain high reward while still approaching a support boundary;
the two outcomes are not treated as equivalent.

\paragraph{Shared use across categorical protocols.}
E6-A fixes the concentration and scans the strength of the local
negative update to identify the transition from the positive-only
ceiling to useful semantic extrapolation and then to performance
reversal. E6-B learns the concentration and applies near/far causal
interventions and matched controls. E6-C compares aligned and shuffled
policy embeddings to determine whether improvements on the hidden action
require the specified semantic correspondence.

E5 serves a distinct but complementary role. Its direct-softmax
diagnostic and categorical causal reconstruction isolate persistent
surprisal growth, logit-gap growth, and simplex-boundary suppression
under repeated negative updates. It shares the fixed-advantage and
near/far intervention logic of D-U1, but it does not use the 64-action
shared-semantic catalogue and does not support the E6 semantic
generalization claim.

\subsection{D4RL Locomotion Datasets}
\label{app:hopper_datasets}

The task-level evaluation uses HalfCheetah, Hopper, and Walker2d under
the medium, medium-replay, and medium-expert datasets
\citep{fu2020d4rl}; the mechanism diagnostics use Hopper, described
below.

Hopper is a continuous-control task with an
11-dimensional observation space and a three-dimensional bounded action
space. We use the D4RL
\texttt{hopper-medium},
\texttt{hopper-medium-replay}, and
\texttt{hopper-medium-expert}
datasets. The medium dataset contains trajectories collected from a
partially trained policy, medium-replay contains the corresponding replay
buffer collected throughout training, and medium-expert combines medium
and expert behavior. Hopper-medium-replay is used for the primary mechanism
diagnostic because its broader behavioral coverage provides a wider range of
learner-relative distances, while all three Hopper datasets are retained for
the task-level performance evaluation.

\subsection{Structured-Generation Tasks and Dataset Construction}
\label{app:structured_generation_dataset}

The structured-generation evaluation contains nine tasks:
Countdown, Word Sorting, Spiral Matrix, Mini Sudoku, Maze, Word Ladder,
Knights \& Knaves, Graph Coloring, and WikiSQL. Countdown follows the
arithmetic task of \citet{hutton2002countdown} and the verifier-guided
LLM formulation used by TinyZero~\citep{pan2025tinyzero}; Word Sorting,
Spiral Matrix, Mini Sudoku, Maze, Word Ladder, Knights \& Knaves, and
Graph Coloring are generated from Reasoning Gym
\citep{stojanovski2025reasoninggym}; WikiSQL follows
\citet{zhong2017seq2sql}. The corresponding source revisions are fixed by
the registered experiment configuration.

All tasks use a common prompt-indexed replay-bank schema containing the
task identity, prompt, verifier-confirmed oracle completion, verifier
outcomes, model-independent negative completions, and task metadata.
Learner-relative quantities are excluded from bank construction:
current-policy probability, surprisal, near/far labels, and taper weights
are recomputed from the learner when required rather than stored with the
data. The replay banks remain frozen during optimization.

\begin{table}[t]
\centering
\caption{Structured-generation task interfaces.}
\label{tab:structured_generation_tasks}
\small
\setlength{\tabcolsep}{3.2pt}
\begin{tabular*}{\columnwidth}{@{\extracolsep{\fill}}lll}
\toprule
Task & Source & Verified output \\
\midrule
Countdown & Generated & Arithmetic expression \\
Word Sorting & Reasoning Gym & Word permutation \\
Spiral Matrix & Reasoning Gym & Integer sequence \\
Mini Sudoku & Reasoning Gym & \(4\times4\) grid \\
Maze & Reasoning Gym & Shortest-path length \\
Word Ladder & Reasoning Gym & Word path \\
Knights \& Knaves & Reasoning Gym & Role assignment \\
Graph Coloring & Reasoning Gym & Vertex-color map \\
WikiSQL & WikiSQL & SQL logical form \\
\bottomrule
\end{tabular*}
\end{table}

Each structured-generation task uses 5,000 training prompts and
500 validation prompts under the unified paper-facing protocol. Reserved test
data are not used for coefficient tuning. The structured-generation optimization
interface uses 16 negative responses per training prompt, keeps the replay data
fixed during optimization, and recomputes learner-relative mean completion-token
surprisal from the current policy. This zero-refresh construction isolates
stale-negative reuse without storing learner-relative remoteness in the bank.

\subsection{Training and Evaluation Protocols}

\section{Controlled Experimental Protocols}
\label{app:controlled_protocols}

This section specifies the interventions used in
Sections~\ref{sec:controlled_identification}
and~\ref{sec:repulsion_control}. The definitions of C-U1 and D-U1 are given in
Appendix~\ref{app:experimental_details}. Network architectures,
optimizer settings, training budgets, seeds, and evaluation frequencies
are reported in the protocol-specific subsections below.

\subsection{Shared Notation and Experimental Invariants}
\label{app:controlled_shared}

For a negative sample
$z=(s,a)$, define its raw repulsive update as
\begin{equation}
    g_\theta^-(z)
    =
    -\widehat A^-(z)
    \nabla_\theta
    \log\pi_\theta(a\mid s),
    \qquad
    \widehat A^-(z)
    =
    \max\{-\widehat A(z),0\}.
    \label{eq:controlled_negative_update}
\end{equation}
Learner-relative remoteness is measured by
\begin{equation}
    D_\theta(z)
    =
    -\log\pi_\theta(a\mid s),
    \label{eq:controlled_remoteness}
\end{equation}
or by its equivalent standardized Gaussian component when additive
density constants are irrelevant.

Unless explicitly stated otherwise, all branches of one controlled
comparison share:

\begin{enumerate}
    \item the same training and evaluation contexts;
    \item the same positive and negative samples;
    \item the same fixed rewards or advantage labels;
    \item the same initial policy parameters;
    \item the same minibatch order, optimizer, and training horizon;
    \item the same evaluation checkpoints and terminal classification
    rules.
\end{enumerate}

Only the variable named by the intervention is changed. Actions generated
under the same context are aggregated at the context level and are not
treated as independent contexts for statistical inference.


\subsection{Source-Isolation Protocol}
\label{app:source_isolation_protocol}

\paragraph{Purpose.}
The source-isolation experiment asks whether policy-relative remoteness
changes negative-update magnitude when sample quality is held fixed. It
identifies the source of a per-sample difference, not its causal effect on
training.

\paragraph{Continuous product-manifold construction.}
For each C-U1 context, we select multiple negative actions from a common
reward contour:
\begin{equation}
    R(s,a_i)
    =
    R(s,a_j),
    \qquad
    \widehat A(s,a_i)
    =
    \widehat A(s,a_j)
    <0,
    \label{eq:source_equal_quality}
\end{equation}
while ensuring that
\begin{equation}
    D_\theta(s,a_i)
    \neq
    D_\theta(s,a_j).
    \label{eq:source_different_remoteness}
\end{equation}
The reward coordinate and the policy-distance coordinate are therefore
independently controlled. This exact counterfactual comparison is the
reason for using a synthetic product-manifold construction: an external
offline dataset generally cannot provide two actions with precisely
matched task quality and independently prescribed distance from the
same current policy.

The comparison is performed both at policy initialization and after
positive-only pretraining. No negative sample used as a probe is applied
as an optimizer update during this protocol. Let \(u_\theta(s)\) denote
the policy-output coordinate used by the controlled probe, e.g.,
fixed-covariance Gaussian mean coordinates in C-U1 and direct logits in
D-U1. For every probe, we record
\begin{equation}
\begin{aligned}
I_{\mathrm{score}}(z)
&=
\left\|
\nabla_{u}\log \pi_{u}(a\mid s)
\right\|_2,\\
I_{\mathrm{sample}}(z)
&=
|\widehat A(z)|\, I_{\mathrm{score}}(z).
\end{aligned}
\label{eq:source_influence_metrics}
\end{equation}
Because \(|\widehat A|\) is identical across the matched actions,
differences in \(I_{\mathrm{sample}}\) arise from policy-score geometry
rather than advantage magnitude.

For neural implementations, we additionally record the full
parameter-gradient norm and the aggregate gradient obtained from samples
in the same remoteness range. These implementation-level actor-gradient
diagnostics are used as transfer checks, not as theorem-level score
measures.

We report near-to-far ratios of remoteness, score norm, single-sample
gradient norm, aggregate gradient norm, and gradient-direction
coherence. Ratios are first computed within each context and seed and
are then aggregated across the prespecified seeds.

\paragraph{Categorical counterpart.}
D-U1 uses unordered categorical actions with identical fixed negative
labels. The comparison varies the current probability assigned to the
negative action while preserving its reward role and advantage:
\begin{equation}
    \widehat A(s,i)
    =
    \widehat A(s,j)
    <0,
    \qquad
    \pi_\theta(i\mid s)
    \neq
    \pi_\theta(j\mid s).
    \label{eq:categorical_source_match}
\end{equation}
For a direct categorical logit parameterization, the score norm is
bounded and need not grow without limit as probability decreases.
The relevant far-field effect is instead persistence: the negative
update continues to increase the selected action's logit gap even after
its probability has become small. We therefore report probability,
surprisal, logit gap, score norm, and cumulative suppression over
repeated updates. This distinguishes categorical support suppression
from the unbounded Gaussian-distance amplification measured in C-U1.


\subsection{Causal-Transmission Protocol}
\label{app:causal_transmission_protocol}

\paragraph{Purpose.}
The source-isolation protocol identifies the source of abnormal update
magnitude but not its effect on task behavior. Causal transmission is tested in the nonlinear
Gaussian protocol, where the policy mean and variance evolve jointly,
and in the categorical support reconstruction of D-U1.

\paragraph{Dynamic near--far partition.}
At each intervention step, negative samples are partitioned using their
current-policy remoteness. Let
$\mathcal N_\theta$
and
$\mathcal F_\theta$
denote the prespecified near- and far-field sets. Their membership is
recomputed as the policy evolves rather than frozen at initialization.
The thresholds or quantiles are fixed before training and specified with the
corresponding environment.

\paragraph{Interventions.}
For a negative sample $z$, the following multipliers are applied to
$g_\theta^-(z)$:
\begin{equation}
\begin{aligned}
    w_{\mathrm{uncontrolled}}(z)
    &=1,\\
    w_{\mathrm{positive}}(z)
    &=0,\\
    w_{\mathrm{near\text{-}zero}}(z)
    &=
    \mathbf 1[z\notin\mathcal N_\theta],\\
    w_{\mathrm{far\text{-}zero}}(z)
    &=
    \mathbf 1[z\notin\mathcal F_\theta].
\end{aligned}
\label{eq:causal_zero_interventions}
\end{equation}

The far-cap intervention preserves the direction of a far-field update
but limits its detached magnitude to a reference level estimated from
near-field negative samples:
\begin{equation}
    w_{\mathrm{far\text{-}cap}}(z)
    =
    \begin{cases}
        \displaystyle
        \min\left\{
            1,
            \frac{C_{\mathrm{near}}}
                 {I_{\mathrm{sample}}(z)+\epsilon}
        \right\},
        & z\in\mathcal F_\theta,\\[3mm]
        1,
        & \text{otherwise},
    \end{cases}
    \label{eq:causal_far_cap}
\end{equation}
where $C_{\mathrm{near}}$ is the prespecified near-field update-magnitude
reference.

To distinguish selective intervention from a generic decrease in
negative optimization, the global-matched control applies one scalar
to every negative sample:
\begin{equation}
    w_{\mathrm{global}}(z)
    =
    \alpha_{\mathrm{match}}.
    \label{eq:causal_global_weight}
\end{equation}
The scalar is chosen using the training batch so that its detached aggregate
negative-update magnitude matches the corresponding selective
intervention:
\begin{equation}
    \alpha_{\mathrm{match}}
    \sum_{z\in\mathcal B^-}
        I_{\mathrm{sample}}(z)
    =
    \sum_{z\in\mathcal B^-}
        w_{\mathrm{selective}}(z)
        I_{\mathrm{sample}}(z).
    \label{eq:causal_budget_match}
\end{equation}
The matching coefficient is not differentiated through.

\paragraph{Causal contrasts.}
The primary causal contrast is between removing near-field and
far-field negative updates. If Near-zero leaves the dynamics unchanged
while Far-zero or Far-cap alters them, the effect cannot be attributed
to negative feedback in general. The global-matched branch quantifies
how much of the change follows from reducing total negative-update magnitude
and how much requires selecting samples by remoteness.

\paragraph{Temporal and terminal measurements.}
We record policy displacement, reward, negative-gradient magnitude,
mean dynamics and, for learnable-variance runs, scale dynamics, entropy or effective support, and the first
time at which each prespecified event occurs. The analysis checks whether
far-field update magnitude rises before policy drift and whether drift precedes
task-performance loss.

Terminal outcomes are classified separately as:

\begin{enumerate}
    \item \textbf{task-performance collapse}: a sustained loss of the
    prespecified task metric while optimization remains finite;
    \item \textbf{support or variance-boundary event}: categorical
    effective support or Gaussian variance crosses its prespecified
    boundary without requiring NaN or Inf;
    \item \textbf{numerical collapse}: a parameter, loss, gradient, or
    model output becomes NaN or Inf.
\end{enumerate}

A boundary event is not automatically counted as task-performance
collapse, and finite execution is not automatically counted as
convergence.

\subsection{Controlled Negative-Strength Sweep}
\label{app:negative_strength_sweep}

The controlled negative-strength sweep varies only the effective
negative-repulsion strength $q/p$, where $p$ denotes the aggregate
positive-attraction budget and $q$ denotes the aggregate
negative-repulsion budget. All sweep arms share the same data, fixed
labels, initialization, optimizer, training horizon, and event
thresholds. The positive-only baseline corresponds to $q/p=0$.

Held-out-context reward evaluates unseen-state generalization within the
C-U1 controlled protocol. Policy shift is the normalized displacement of
the learned policy from the positive-only target and visualizes how far
negative feedback pushes the policy beyond the imitation solution.
Task-performance collapse, support or variance-boundary events, and
NaN/Inf numerical failures are reported as separate event types.


\subsection{Controlled Taper Comparison}
\label{app:controlled_taper_protocol}

\paragraph{Controlled utility construction.}
In C-U1, the hidden optimum $a_\star(s)$ lies beyond the center of the
positive demonstrations $a_+(s)$. A local negative anchor is placed on
the opposite side of $a_+(s)$, so repulsion from this anchor points
toward $a_\star(s)$. Additional negative actions are placed on matched
reward contours at larger policy-relative distances. Increasing
negative strength therefore produces three observable regimes:
positive-only under-extrapolation, useful local extrapolation, and
far-field over-extrapolation or instability.

D-U1 implements the same logic with unordered semantic actions. The
hidden optimal target lies beyond the demonstrated positive direction,
while the local negative target lies in the opposite semantic
direction. Repelling the local negative action can therefore improve
the shared representation toward the hidden optimal action. Additional
remote negative actions create support pressure without changing their
fixed negative label. Evaluation on independently sampled contexts is
reported as held-out-context generalization.

Let
\[
x(D)
=
\left[
\frac{D-\tau}{c}
\right]_{+},
\qquad
r_x(D)=\sqrt{x(D)}.
\]
The Linear and Quadratic names refer to polynomial order in the
radialized coordinate \(r_\theta=\sqrt{x_\theta}\), rather than in the
squared-remoteness coordinate \(x_\theta\).
\paragraph{Weighting rules.}
For negative-sample remoteness $D$, the compared rules are
\begin{equation}
\begin{aligned}
    \omega_{\mathrm{positive}}(D)
    &=0,\\
    \omega_{\mathrm{uncontrolled}}(D)
    &=1,\\
    \omega_{\mathrm{global}}(D)
    &=\alpha,\\
    \omega_{\mathrm{hard}}(D)
    &=
    \mathbf 1[D\leq\tau],\\
    \omega_{\mathrm{Rec\text{-}L}}(D)
    &=
    \frac{1}{1+\lambda r_x(D)},\\
    \omega_{\mathrm{Rec\text{-}Q}}(D)
    &=
    \frac{1}{1+\lambda r_x(D)^2},\\
    \omega_{\mathrm{exp}}(D)
    &=
    \exp\{-\lambda r_x(D)^2\}.
\end{aligned}
\label{eq:controlled_taper_rules}
\end{equation}
All weighting values are computed from detached remoteness statistics;
the policy cannot reduce its loss by differentiating through the taper.

\paragraph{Near-field-retention matching.}
Let $\mathcal N$ denote the prespecified useful near-field set and define
its retained first-order update magnitude as
\begin{equation}
    \rho_{\mathrm{near}}(\omega)
    =
    \frac{
        \sum_{z\in\mathcal N}
        \omega(D_\theta(z))
        I_{\mathrm{sample}}(z)
    }{
        \sum_{z\in\mathcal N}
        I_{\mathrm{sample}}(z)
    }.
    \label{eq:near_retention}
\end{equation}
Taper parameters are calibrated on training contexts so that the
compared selective rules retain the same target near-field fraction,
up to the predetermined tolerance. Parameters are then frozen before
evaluation. This comparison isolates how the rules treat the far-field
tail after preserving comparable local utility.

\paragraph{Aggregate-budget matching.}
For a weighting function $\omega$, define its detached negative-update
budget as
\begin{equation}
    B(\omega)
    =
    \sum_{z\in\mathcal B^-}
    \omega(D_\theta(z))
    I_{\mathrm{sample}}(z).
    \label{eq:taper_budget}
\end{equation}
The budget-matched global coefficient is
\begin{equation}
    \alpha_{\mathrm{budget}}
    =
    \frac{
        B(\omega)
    }{
        B(\omega_{\mathrm{uncontrolled}})
    }.
    \label{eq:taper_global_match}
\end{equation}
This control applies the same total first-order negative-update budget
without using remoteness. A selective taper can therefore be credited
only for improvements that remain after comparison with its
budget-matched global baseline.

\paragraph{Measurements.}
Each method is evaluated over the full prespecified training horizon.
We report task performance, held-out-context generalization, distance
from the hidden optimum, policy displacement, entropy or variance,
effective support, near-field retention, aggregate negative-update
budget, and all three terminal-event categories. Method comparisons use
the same seeds and paired statistics.

\begin{figure}[t]
\centering
\includegraphics[
    width=\columnwidth
]{figures/fig_6_4_2_leftfig_bigtext_legend_protocol.pdf}
\caption{
\textbf{Controlled taper-shape comparison.}
Matched operating points (left: near-field-retention matched; right: matched
for \(B(\omega)\), the weighted sum of per-sample output-coordinate update
magnitudes). Bars split \(B(\omega)\) into near- and far-field contributions;
dashed traces show held-out-context reward.
}
\end{figure}

We compare positive-only learning, global negative scaling,
reciprocal-linear tapering, reciprocal-quadratic tapering, and
exponential tapering while changing only the negative weighting rule. Two
matched comparisons are used. Near-field-retention matching preserves
comparable local negative-update magnitude and isolates differences in far-field
decay, whereas aggregate-budget matching tests whether a selective taper
outperforms a global reduction with the same total negative-update
budget.

Under both matched comparisons, exponential tapering retains near-field
repulsion while suppressing the far-field tail more rapidly than polynomial
tapers. Positive-only removes both components, whereas global scaling reduces
them indiscriminately.



\subsection{External Taper Validation}
\label{app:external_taper_protocol}

\paragraph{Scientific scope.}
External taper experiments test whether the controlled weighting
principle transfers to complex policies. They do not reproduce the exact
counterfactual geometry of C-U1 or D-U1 and therefore do not replace the
controlled source and causal experiments.

\paragraph{Hopper.}
All Hopper taper methods branch from the same actor initialization and
use the same offline transitions and advantage labels. For an action
with inverse-squash coordinate $u$, Gaussian mean $\mu_\theta(s)$, and
diagonal scale $\sigma_\theta(s)$, the remoteness statistic is
\begin{equation}
    D_{\mathrm{Hopper}}(s,a)
    =
    \frac{1}{2}
    \sum_j
    \left(
        \frac{
            u_j-\mu_{\theta,j}(s)
        }{
            \sigma_{\theta,j}(s)
        }
    \right)^2.
    \label{eq:hopper_taper_remoteness}
\end{equation}
The additive Gaussian normalization term is omitted because it does not
change the ordering used by the taper. Remoteness is recomputed under
the current actor and detached before constructing the negative weight.
The mechanism protocol with frozen advantages is kept separate from the
standard offline-RL performance protocol; the former studies transfer of
the far-field control mechanism, whereas the latter reports final
environment return.

\paragraph{Countdown.}
For each replay prompt, the V2 record contains a verified oracle
solution and 9--16 unique programmatically generated negative
expressions stratified by construction-time error type. A fixed-width
representation with 16 negative slots was also materialized for
downstream compatibility: 4{,}943 prompts already contained 16 unique
negatives, while 1{,}057 prompts were padded by cycling existing exact
expressions. This representation records conversion provenance rather
than the multiplicity used by the paper-facing objectives. The
paper-facing trainers deduplicate exact expressions and optimize over
all unique negatives available for each prompt. Learner-relative
near/far status is never stored and is recomputed only by methods that
use the corresponding remoteness coordinate.

All protocol-matched V2 Countdown methods use the same prompt set,
verified oracle solutions, frozen V2 negative bank, validation set,
and maximum training budget. Initialization, sampling, calibration,
and ratio-denominator rules are method-specific and are stated
separately.

\paragraph{Comparisons and reporting.}
The common external comparisons are Positive-only, Uncontrolled,
Global scaling, and exponential DRPO. Domain-specific baselines are
reported in the overall-performance section rather than treated as
controlled taper variants. Hopper reports normalized return and policy
statistics; Countdown reports greedy verifier success, Pass@$k$, valid
expression rate, surprisal, and support statistics. Task-performance,
boundary, and NaN/Inf outcomes remain separate in both domains.
\subsection{Shared Objectives and Negative-Update Controls}
\label{app:shared_methods}

This section defines the protocol controls and remoteness-aware
weighting rules shared by the D4RL and structured-generation experiments.
Their task-specific remoteness signals, data construction, and optimization
protocols are described separately in
Appendices~\ref{app:d4rl_methods} and
\ref{app:structured_generation_methods}.

\paragraph{Shared signed objective.}

Let \(z\) denote an offline transition or response, let
\(A(z)\) denote its fixed signed learning signal during an actor update,
and let \(D_\theta(z)\) denote its learner-relative remoteness under the
current policy. We write

\[
A^{+}(z)=\max\{A(z),0\},
\qquad
A^{-}(z)=\max\{-A(z),0\}.
\]

The shared family of actor objectives is

\[
\mathcal{J}_{\omega}(\theta)
=
\mathbb{E}_{z\sim\mathcal{B}}
\left[
A^{+}(z)\log \pi_\theta(z)
 -
\omega\!\left(D_\theta(z)\right)
A^{-}(z)\log \pi_\theta(z)
\right],
\]

where positive updates retain unit weight and
\(\omega(D)\) controls only the negative branch. During each actor
update, the advantage and remoteness-derived weight are treated as
fixed coefficients; gradients are not propagated through the weighting
rule.

For the remoteness-aware methods, define the normalized excess
remoteness

\[
x_\theta(z)
=
\left[
\frac{D_\theta(z)-\tau}{c}
\right]_{+},
\]

where \(\tau\) is the near-field threshold, \(c>0\) is the task-specific
scale, and \([u]_{+}=\max\{u,0\}\). Consequently, samples inside the
near-field region \(D_\theta(z)\leq\tau\) retain their full negative
weight.

\paragraph{Positive-only.}

Positive-only removes all negative updates:

\[
\omega_{\mathrm{pos}}(D)=0.
\]

It therefore provides an extreme no-repulsion reference. In D4RL it
retains only transitions with positive actor advantages; in
structured-generation tasks it retains only verifier-confirmed positive
responses under the corresponding task signal.

\paragraph{Global-\(\alpha\).}

Global-\(\alpha\) applies the same attenuation to every negative sample:

\[
\omega_{\mathrm{global}}(D)=\alpha,
\qquad
0\leq\alpha\leq 1.
\]

This control tests whether performance improvements can be explained by
a global reduction in negative-update magnitude, without using
learner-relative remoteness.

\paragraph{Reciprocal-Linear and Reciprocal-Quadratic.}
For the remoteness-aware methods, define the normalized excess
remoteness
\[
x_\theta(z)
=
\left[
\frac{D_\theta(z)-\tau}{c}
\right]_{+},
\]
and its radialized coordinate
\[
r_\theta(z)=\sqrt{x_\theta(z)}.
\]
For Gaussian policies, \(r_\theta\) is asymptotically proportional to
standardized distance in the far field; for sequence policies, it is a
surprisal-derived radial coordinate rather than a geometric distance in
the discrete action space.
The two reciprocal controls use the same remoteness signal and
near-field threshold as DRPO, but differ in their far-field decay:

\[
\omega_{\mathrm{Rec\text{-}L}}(D)
=
\frac{1}{1+\lambda r_\theta(z)},
\]

and

\[
\omega_{\mathrm{Rec\text{-}Q}}(D)
=
\frac{1}{1+\lambda r_\theta(z)^2}
=
\frac{1}{1+\lambda x_\theta(z)}.
\]

Reciprocal-Linear has the slowest far-field attenuation. The quadratic
variant decays more rapidly but retains a polynomial tail. These
controls isolate whether the benefit comes merely from using a
remoteness signal or specifically from the decay profile selected by
DRPO.

\paragraph{DRPO.}

DRPO uses an exponential remoteness taper:

\[
\omega_{\mathrm{DRPO}}(D)
=
\exp\left\{
-\lambda x_\theta(z)
\right\}.
\]

The weight equals one throughout the near-field region and then decays
exponentially once remoteness exceeds \(\tau\). Thus, DRPO preserves
local negative feedback while suppressing the remote negative tail more
rapidly than the reciprocal alternatives.

\paragraph{Hyperparameter selection.}

Within each task, \(\tau\), \(c\), \(\lambda\), and the
Global-\(\alpha\) coefficient are selected using the registered
validation protocol only. Test-task performance is not used for
hyperparameter selection.


\subsection{D4RL Baselines and Experimental Protocol}
\label{app:d4rl_methods}

The definitions of Positive-only, Global-\(\alpha\),
Reciprocal-Linear, Reciprocal-Quadratic, and DRPO follow
Appendix~\ref{app:shared_methods}. This section specifies only the
D4RL-specific baselines, remoteness construction, training protocol,
and evaluation procedure.

\paragraph{Tasks and datasets.}

We evaluate HalfCheetah, Hopper, and Walker2d on the medium,
medium-replay, and medium-expert datasets, producing the nine dataset cells
reported in Table~\ref{tab:d4rl_performance}. The same fixed offline
dataset is used by every protocol-matched method within a task.

\paragraph{Published offline-RL baselines.}

CQL, TD3+BC, and IQL are included as published reference results.
These values are not protocol-matched reruns and are therefore marked
by \(\dagger\) in Table~\ref{tab:d4rl_performance}. The main table
reports their published task-level means. Their original uncertainty
statistics and exact score provenance are recorded separately from the
protocol-matched DRPO experiments.

\paragraph{Gaussian-policy remoteness.}

For an action \(a\in\mathbb{R}^{d_a}\) and current Gaussian policy with
mean \(\mu_\theta(s)\) and standard deviation \(\sigma_\theta(s)\), we use the
following implementation-level standardized remoteness statistic

\[
D_{\mathrm{D4RL}}(s,a)
=
\frac{1}{d_a}
\sum_{j=1}^{d_a}
\left(
\frac{a_j-\mu_{\theta,j}(s)}
     {\sigma_{\theta,j}(s)}
\right)^2.
\]

This dimension-normalized convention differs from
Equation~\eqref{eq:hopper_taper_remoteness} only by a fixed task-level scale.
The associated threshold and scale parameters absorb that constant, and each
protocol keeps its convention fixed across all compared methods.

When the actor uses a transformed action distribution, this quantity is
computed in the policy's native Gaussian coordinates. The resulting
remoteness is used only to determine the negative-update coefficient;
the distance term is detached during the actor update.

\paragraph{Protocol-matched actor comparison.}

Positive-only, Global-\(\alpha\), Reciprocal-Linear,
Reciprocal-Quadratic, and DRPO share the same offline dataset, critic
checkpoint or critic-training protocol, actor initialization, network
architecture, optimizer, minibatch schedule, actor-update budget,
evaluation schedule, and random seeds. They differ only in the
negative-update weighting rule defined in
Appendix~\ref{app:shared_methods}.

The common actor advantage construction is held fixed across these
methods. No method is allowed to change the critic, relabel samples, or
alter the data distribution in order to improve its actor result.

\paragraph{Hyperparameters.}

The registered training and selection configuration is summarized in
Table~\ref{tab:d4rl_registered_config}. Although the remoteness
coordinate is defined consistently across tasks, its empirical
distribution depends on the dataset. We therefore register
task-specific development grids so that each controller covers
comparable operating regimes under the corresponding remoteness
distribution. All grids and selection rules are frozen before
held-out evaluation and are populated directly from the experiment
configuration, rather than reconstructed or adjusted after observing
test results.

\begin{table}[t]
\centering
\caption{Configuration of the D4RL-9 development sweeps. Each
task-specific grid contains ten candidates: five from the coarse sweep
and five from refinement.}
\label{tab:d4rl_registered_config}

\small
\setlength{\tabcolsep}{6pt}
\renewcommand{\arraystretch}{1.10}

\begin{tabular}{@{}ll@{}}
\toprule
Configuration & Value \\
\midrule
Policy-update budget
& $1{,}000{,}000$ \\
Evaluation
& Every $50{,}000$ updates; $10$ episodes \\
Fixed geometry
& $\tau=0$ \\
Global scaling, $\alpha$
& Task-specific grid \\
Reciprocal coefficients, $c$
& Task-specific grids for Linear and Quadratic \\
Exponential multiplier
& $\alpha_{\mathrm{Exp}}=1$ \\
\bottomrule
\end{tabular}
\end{table}

\paragraph{Evaluation and uncertainty.}

Task performance is reported using the standard D4RL normalized return.
For every protocol-matched method, the main table reports the mean over
the registered seeds. The D4RL-9 total is the sum of
the nine task-level means.

Task-performance collapse is assessed from normalized return and
training trajectories. Variance or support-boundary events are reported
separately from task performance. NaN/Inf numerical failures are
reported as a third outcome and are not merged with either low return or
policy-boundary events. Any claim concerning convergence, collapse, or
method ranking is based on the registered terminal audit rather than an
intermediate checkpoint.


\subsection{Structured-Generation Baselines and Experimental Protocol}
\label{app:structured_generation_methods}

All structured-generation experiments use
Qwen2.5-0.5B-Instruct~\citep{qwen2024qwen25} with LoRA adaptation and
frozen prompt-indexed replay data. Learner-relative remoteness for DRPO is
the current-policy mean completion-token surprisal and is recomputed during
training rather than stored in the replay bank. Protocol-matched methods
share the registered task data, verifier labels, optimization budget, and
evaluation schedule within each task.

\paragraph{AsymRE.}

Asymmetric REINFORCE defines the response-level advantage using a
tunable reward baseline \(V\):

\[
A_{\mathrm{AsymRE}}(x,y)
=
r(x,y)-V,
\]

and optimizes

\[
\mathcal{J}_{\mathrm{AsymRE}}(\theta)
=
\mathbb{E}_{(x,y)\sim\mathcal{B}}
\left[
\bigl(r(x,y)-V\bigr)
\log\pi_\theta(y\mid x)
\right].
\]

For a binary verifier reward, lowering \(V\) emphasizes successful
responses, whereas increasing \(V\) strengthens suppression of failed
responses. In our branch-balanced signed-reward implementation,
\(r\in\{-1,+1\}\) gives the empirical baseline
\(\widehat V=0\). The swept offset therefore determines the
objective baseline as
\(V=\widehat V+\delta_v=\delta_v\). The baseline is selected using the validation protocol and
is held fixed within each run.

Global-\(\alpha\) and AsymRE remain separate comparisons because they
operate on different underlying learning signals. Global-\(\alpha\)
scales the negative part of the common precomputed signed advantage,
whereas AsymRE constructs its signed advantage directly from the raw
verifier reward and the published baseline parameterization. Under
binary rewards and simplified normalization they are closely related,
but the formal comparison preserves their native objectives and tuning
rules.

\paragraph{\(\beta\)-TOPR (joint fitted reference).}

Let \(\mathcal X\) be the training-prompt set. For each
\(x\in\mathcal X\), let \(y_x^+\) be the verified oracle solution and
let \(\mathcal N_x\) be the set of unique negative completions after
canonical-completion deduplication. Define the mean completion-token log
probability

\[
\bar{\ell}_\theta(x,y)
=
\frac{1}{|y|}
\sum_{t=1}^{|y|}
\log\pi_\theta(y_t\mid x,y_{<t}),
\]

and the full-completion summed log probability used only for the ratio,

\[
s_\theta(x,y)
=
|y|\,\bar{\ell}_\theta(x,y)
=
\sum_{t=1}^{|y|}
\log\pi_\theta(y_t\mid x,y_{<t}).
\]

For the jointly fitted reference adapter \(\mu_{\mathrm{ref}}\), define

\[
s_{\mathrm{ref}}(x,y)
=
\sum_{t=1}^{|y|}
\log\mu_{\mathrm{ref}}(y_t\mid x,y_{<t}),
\]

\[
\rho_\theta(x,y)
=
\exp\!\left(s_\theta(x,y)-s_{\mathrm{ref}}(x,y)\right),
\]

and

\[
\begin{aligned}
w_\beta(x,y)
&=
\exp\!\left(
\beta\min\{s_\theta(x,y)-s_{\mathrm{ref}}(x,y),0\}
\right) \\
&=
\min\{\rho_\theta(x,y)^\beta,1\}.
\end{aligned}
\]

The policy objective is

\[
\begin{aligned}
\mathcal J_{\beta\text{-}\mathrm{TOPR}}(\theta)
={}&
\mathbb E_{x\sim\operatorname{Unif}(\mathcal X)}
\Bigg[
\bar{\ell}_\theta(x,y_x^+) \\
&-
\frac{1}{|\mathcal N_x|}
\sum_{y\in\mathcal N_x}
\operatorname{sg}[w_\beta(x,y)]
\bar{\ell}_\theta(x,y)
\Bigg].
\end{aligned}
\]

with actor update

\[
g_{\beta\text{-}\mathrm{TOPR}}(\theta)
=
\nabla_\theta\mathcal J_{\beta\text{-}\mathrm{TOPR}}(\theta).
\]

This objective samples prompts uniformly, assigns unit total
coefficient to the positive branch and unit total coefficient to the
negative branch, and averages unique negatives within each prompt. It
is not a simple uniform average over all response rows. The
ratio-derived weight is detached from the policy gradient.

The reference adapter is fitted with

\[
\mathcal J_{\mathrm{ref}}
=
\mathbb E_{x\sim\operatorname{Unif}(\mathcal X)}
\left[
\frac{1}{2}\bar{\ell}_{\mathrm{ref}}(x,y_x^+)
+
\frac{1}{2|\mathcal N_x|}
\sum_{y\in\mathcal N_x}
\bar{\ell}_{\mathrm{ref}}(x,y)
\right].
\]

The reference adapter is updated once per policy step. Reference
outputs and ratio-derived weights are detached from the policy
gradient. The 0.5/0.5 masses belong only to the reference-fitting
objective, not to the policy objective. The data records do not contain
behavior-policy log probabilities; the ratio denominator is supplied by
the jointly fitted reference adapter during training. The completed
TOPR coefficient scan is a provenance-limited development pilot and does
not establish convergence, statistical significance, steady state, or
a formal method ranking.

\paragraph{DPO.}

Canonical DPO~\citep{rafailov2023direct} is trained with the same LoRA
parameterization as the other structured-generation methods. Its policy and
frozen reference are initialized from the registered short-SFT checkpoint,
rather than directly from the pretrained backbone; this difference is marked
by \(^{\ddagger}\) in Table~\ref{tab:structured_generation_performance}.
For each prompt \(x\), the verified oracle solution
\(y_x^+\) is paired with every unique verifier-incorrect completion
\(y\in\mathcal N_x\). Using the summed completion log probability
\(s_\theta\) defined above, the loss is

\[
\begin{aligned}
\mathcal L_{\mathrm{DPO}}(\theta)
={}&-
\mathbb E_{x\sim\operatorname{Unif}(\mathcal X)}
\Bigg[
\frac{1}{|\mathcal N_x|}
\sum_{y\in\mathcal N_x}
\log\sigma\Bigl(
\beta\bigl[
 s_\theta(x,y_x^+)-s_{\mathrm{ref}}(x,y_x^+) \\
&\hspace{8.2em}
-s_\theta(x,y)+s_{\mathrm{ref}}(x,y)
\bigr]
\Bigr)
\Bigg].
\end{aligned}
\]

The loss is averaged over unique negatives within each prompt and then
uniformly across prompts. The reference remains frozen throughout the
fixed 1,200-step continuation, and \(\beta\) is selected using the
registered validation sweep.

\paragraph{Method-specific remoteness.}

For DRPO and the other remoteness-taper controls, the learner-relative
coordinate for a prompt \(x\) and completion
\(y=(y_1,\ldots,y_T)\) is the current mean completion-token surprisal

\[
D_{\mathrm{SG}}(x,y)
=
-\frac{1}{T}
\sum_{t=1}^{T}
\log
\pi_\theta
\left(
y_t\mid x,y_{<t}
\right).
\]

Only completion tokens are included; prompt tokens and padding positions
are excluded. Mean rather than summed surprisal is used so that response
length does not mechanically determine DRPO remoteness. The remoteness
score is detached before applying the shared taper functions. AsymRE
does not use remoteness and applies its signed reward directly. The
Joint Fitted-Reference \(\beta\)-TOPR variant uses the dynamic
full-completion policy/reference log-ratio defined above. Canonical DPO
instead optimizes the frozen-reference pairwise log-ratio in its own
objective; neither baseline uses DRPO's mean-token remoteness.

\paragraph{Frozen response-bank construction.}

Before policy optimization, the registered model-independent generator
constructs a frozen prompt-indexed replay bank for each
structured-generation task. Each training prompt has one verifier-confirmed
oracle completion and 16 model-independent verifier-incorrect completions
under the registered structured-generation training interface. The
paper-facing trainers canonicalize exact outputs before consumption. The
bank stores no learner-relative near/far label and no behavior-policy log
probability. Verifier outcome, response length, output-validity status, and
task-specific error metadata remain separate stored diagnostics even when
they share the same binary training reward.

\paragraph{Training and checkpoint selection.}

All protocol-matched methods use Qwen2.5-0.5B-Instruct
\citep{qwen2024qwen25} with the registered LoRA parameterization. The
Joint Fitted-Reference \(\beta\)-TOPR arm additionally creates a fresh
reference LoRA with parameters copied from the initial policy adapter.
Within each structured-generation task, protocol-matched methods share the
AdamW optimizer~\citep{loshchilov2019decoupled}, learning-rate schedule,
minibatch construction, optimizer-step budget, validation protocol, and
evaluation schedule unless a method-specific protocol states otherwise.
Hyperparameters are selected using the registered validation rule.

The structured-generation experiments use 5,000 training prompts and
500 validation prompts per task; reserved test data are not used for
coefficient tuning. All tasks use 1,200 optimizer updates, 16 negatives per
training prompt, no early stopping, and Pass@8 as the primary sampled-decoding
metric. The common LoRA configuration is rank 32, alpha 64, and dropout 0.05.
Task-specific runtime overrides such as sequence length and evaluation batch
size follow the registered task adapters and are not varied as scientific
comparison factors.

\paragraph{Verifier-based evaluation.}

Greedy verifier success is the fraction of held-out instances solved by
greedy decoding. Pass@$k$ is the fraction solved by at least one of
the registered \(k\) sampled responses. Task-specific output-validity
diagnostics are reported separately from verifier success where applicable.

A higher verifier success rate is not attributed to improved task
performance when it is accompanied by a deterioration in the corresponding
output-validity diagnostic.

Task-performance degradation, token- or sequence-support boundary
events, and NaN/Inf numerical failures are audited independently.
Claims about stability or final method ranking use the registered
terminal checkpoint audit rather than finite-step pilot behavior.


\paragraph{Baseline coefficient selection.}
For each structured-generation task, AsymRE scans
\(\delta_v\in\{-1.0,-0.9,-0.75,-0.5,0.0\}\), Joint Fitted-Reference
\(\beta\)-TOPR scans \(\beta\in\{0.25,0.5,1.0\}\), and canonical DPO
scans \(\beta\in\{0.05,0.1,0.2\}\). Coefficients are selected on the
registered validation metric using the same task-level selection rule as the
main structured-generation comparison; test data are not used for coefficient
selection.

\subsection{Structured-Generation DRPO Coefficient Response}
\label{app:structured_generation_taper_sensitivity}

Figure~\ref{fig:app_structured_generation_coefficient_response} reports the
DRPO coefficient-response curves across the structured-generation tasks.
Each panel plots late-window Pass@8 against the registered taper coefficient
\(\lambda\), with the corresponding Positive-only result shown as a
horizontal reference.

Across tasks, the response is task dependent rather than a monotone
``larger is better'' relationship. Countdown, Word Ladder, and
Knights \& Knaves show substantial gains over Positive-only over part of the
evaluated coefficient range, while Mini Sudoku and Maze show smaller
improvements. Spiral Matrix is already near the task ceiling. Word Sorting
and WikiSQL improve as \(\lambda\) increases but the observed sweeps do not
establish a closed interior optimum. For Graph Coloring, small
\(\lambda\) is insufficient, an intermediate range is effective, and the
benefit declines as \(\lambda\) becomes large, with high-\(\lambda\)
performance returning toward the Positive-only level. These fixed-horizon
curves characterize coefficient sensitivity and task-dependent useful
operating ranges; they do not establish convergence, steady state, or a
universal cross-task value of \(\lambda\).

\begin{figure*}[t]
    \centering
    \includegraphics[
        width=\textwidth
    ]{figures/fig_app_structured9_drpo_coefficient_response.pdf}
    \caption{
        \textbf{DRPO coefficient response across structured-generation tasks.}
        Each panel reports late-window Pass@8 versus the registered taper
        coefficient \(\lambda\); dashed horizontal lines denote the
        corresponding Positive-only references. Task-specific vertical scales
        are used to make within-task response shape visible.
    }
    \Description{Nine-panel DRPO coefficient-response figure covering
    Countdown, Word Sorting, Spiral Matrix, Mini Sudoku, Maze, Word Ladder,
    Knights and Knaves, Graph Coloring, and WikiSQL.}
    \label{fig:app_structured_generation_coefficient_response}
\end{figure*}

\clearpage
\end{document}
